%======================================================================
%  A Derivative-Free Trust-Funnel Method for Relaxable Simulation-Based
%  Inequality Constraints
%
%  Theory document (POptUS / IBCDFO, branch experimental/relaxable_constraints).
%  Milestone status:
%    M1  Scoping note ......... notation.tex + positioning.tex        [done]
%    M2  Algorithm spec ....... algorithm.tex                         [done]
%    M3  Convergence analysis . convergence.tex                       [done]
%
%  Build:  pdflatex main && bibtex main && pdflatex main && pdflatex main
%======================================================================
\documentclass[11pt]{article}

\usepackage[margin=1in]{geometry}
\usepackage{amsmath,amssymb,amsthm}
\usepackage{booktabs}
\usepackage{url}
\providecommand{\path}[1]{\texttt{#1}}
\usepackage[round,authoryear]{natbib}
\usepackage[colorlinks=true,allcolors=blue]{hyperref}

% --- notation macros ---
\newcommand{\R}{\mathbb{R}}
\newcommand{\Rn}{\mathbb{R}^{n}}
\newcommand{\feas}{\mathcal{F}}
\newcommand{\Lset}{\mathcal{L}}
\newcommand{\A}{\mathcal{A}}
\DeclareMathOperator*{\argmin}{arg\,min}
\DeclareMathOperator{\pred}{pred}
\DeclareMathOperator{\ared}{ared}
\newcommand{\eps}{\varepsilon}

% --- theorem environments ---
\theoremstyle{definition}
\newtheorem{definition}{Definition}
\newtheorem{assumption}{Assumption}
\theoremstyle{plain}
\newtheorem{theorem}{Theorem}
\newtheorem{lemma}{Lemma}
\newtheorem{proposition}{Proposition}
\newtheorem{corollary}{Corollary}
\newtheorem{remark}{Remark}

\title{A Derivative-Free Trust-Funnel Method for\\
       Relaxable Simulation-Based Inequality Constraints}
\author{POptUS / IBCDFO project\\
        \small branch \texttt{experimental/relaxable\_constraints}}
\date{Working draft --- M1--M4 complete; smooth measure; self-contained convergence (TD1+TD2 resolved)}

\begin{document}
\maketitle

\begin{abstract}
We study derivative-free optimization of a simulation-based objective subject to
\emph{relaxable, quantifiable} simulation-based inequality constraints
$c(x)\le0$ --- constraints that return a meaningful violation magnitude at
infeasible points. We propose a trust-funnel (\emph{not} penalty, \emph{not}
feasible-iterate) method that builds fully-linear models of both the objective
and the constraints on every iteration, and we target convergence to a KKT
stationary point under the Mangasarian--Fromovitz constraint qualification. This
document fixes the problem class, assumptions, and notation (M1); states the
algorithm (M2); develops the convergence analysis (M3); and reports an adversarial
self-review with a code-bridge appendix (M4). Using the \emph{smooth} infeasibility
measure $\theta=\tfrac12\|[c]_+\|^2$, we prove that every limit point is feasible
($\theta_k\to0$, Theorem~\ref{thm:T1}) and that, under MFCQ, every feasible limit
point satisfies the KKT conditions ($\lim_k\pi_k^f=0$, Theorem~\ref{thm:T2}). A
central conceptual point is that building fully-linear models on every iteration is
\emph{necessary but not sufficient}: a criticality step and the funnel are the two
further ingredients the proof requires. The analysis is \emph{self-contained} --- the
standard single-region trust-region results it uses are reproduced from first
principles and applied to both the objective and the infeasibility merit --- so the
only inputs are the stated assumptions.
\end{abstract}

\tableofcontents
\bigskip

% ---- M1 content ----
% !TEX root = main.tex
%======================================================================
%  M1 Scoping note --- Problem, assumptions, models, criticality, funnel
%======================================================================
\section{Problem, assumptions, and notation}
\label{sec:problem}

\subsection{The problem class}

We study the inequality-constrained optimization problem
\begin{equation}
  \label{eq:problem}
  \min_{x \in \Rn} \; f(x)
  \quad\text{subject to}\quad
  c(x) \le 0,
\end{equation}
where $f:\Rn\to\R$ and $c=(c_1,\dots,c_p):\Rn\to\R^p$. Both $f$ and $c$ are
available only through a \emph{blackbox / simulation} that returns their values
at a queried point $x$; no derivatives are available and finite differencing is
assumed unreliable or too expensive. We write the feasible set as
\begin{equation}
  \label{eq:feasset}
  \feas \;=\; \{\, x\in\Rn : c(x)\le 0 \,\}.
\end{equation}

The defining structural assumption --- and the source of this project's novelty
--- concerns \emph{what the simulation returns at infeasible points}.

\begin{assumption}[Relaxable, quantifiable, simulation constraints]
  \label{ass:relaxable}
  Each constraint $c_i$ is
  \emph{relaxable} (the simulation can be evaluated at points with $c_i(x)>0$,
  and such evaluations are permitted by the algorithm),
  \emph{quantifiable} (the returned value $c_i(x)$ measures the \emph{degree} of
  satisfaction/violation, not merely a yes/no feasibility flag), and
  \emph{simulation-based} (revealed only by running the blackbox).
  This is the QRSK class of \citet{LeDigabelWild2015}.
\end{assumption}

Relaxability is exactly what \emph{permits} the normal step and the funnel
mechanism of Section~\ref{sec:funnel} to evaluate models built from infeasible
sample points; it is the property that fails for the ``unrelaxable /
feasible-iterate'' methods discussed in Section~\ref{sec:positioning}.

\subsection{Smoothness and the working region}

\begin{assumption}[Smoothness]
  \label{ass:smooth}
  $f$ and every $c_i$ are continuously differentiable on $\Rn$, and their
  gradients $\nabla f$, $\nabla c_i$ are Lipschitz continuous on an open set
  containing the level set $\Lset$ defined below. (These properties are used only
  in the analysis; the \emph{algorithm} never accesses a derivative.)
\end{assumption}

\begin{assumption}[Bounded level set / working region]
  \label{ass:levelset}
  There is a point $x_0$ and a constant such that the region
  \[
    \Lset \;=\; \bigl\{\, x\in\Rn : f(x)\le f(x_0),\ \theta(x)\le\theta_0^{\max}\,\bigr\}
  \]
  swept by the iterates (with $\theta$ and $\theta_0^{\max}$ from
  Sections~\ref{sec:infeas}--\ref{sec:funnel}) is bounded. Consequently $f$, $c$
  and their gradients are bounded on $\Lset$.
\end{assumption}

\begin{assumption}[Feasible, reachable start]
  \label{ass:feasstart}
  A point $x_0$ with $c(x_0)\le 0$ is given (so $\feas\ne\varnothing$).
\end{assumption}

Assumption~\ref{ass:feasstart} is a deliberate scoping decision. Because a
feasible start is available, we need \emph{not} prove global feasibility recovery
from an arbitrary infeasible start, and we drop the usual ``converge to a
stationary point of the infeasibility measure'' fallback dichotomy. The method
may still take \emph{relaxable} excursions into $\{c(x)\not\le0\}$ between
iterates; the funnel of Section~\ref{sec:funnel} drives the infeasibility back to
zero. It also lets the funnel be initialized cleanly (e.g.\ $\theta_0^{\max}$
from the first tolerated excursion).

\subsection{Standing notation}

We use $\|\cdot\|$ for the Euclidean norm unless a subscript indicates otherwise,
$B(x,\Delta)=\{y:\|y-x\|\le\Delta\}$ for the closed trust region, and
$[\,a\,]_+=\max(0,a)$ componentwise. At iteration $k$ the incumbent is $x_k$ with
trust-region radius $\Delta_k>0$. The names of the algorithmic constants are
chosen to match the existing MATLAB scaffold \path{DFO_quad_TR/Algorithm_11_1.m}
and \path{Algorithm_10_2.m}, so the later implementation is a direct
translation:
\[
  \Delta_{\max},\ \eta_0,\ \eta_1,\ \gamma\in(0,1),\ \gamma_{\text{inc}}>1,\
  \epsilon_c,\ \mu,\ \beta,\ \omega\in(0,1),\ r,\ \Delta_{\min}.
\]

\section{Models: fully linear for objective \emph{and} constraints}
\label{sec:models}

At $x_k$ we build a model of the objective, $m_f^k$, and a model of each
constraint, $m_{c,i}^k$, valid on $B(x_k,\Delta_k)$. We adopt the standard DFO
convention that models \emph{interpolate at the incumbent}: $m_f^k(x_k)=f(x_k)$ and
$m_{c,i}^k(x_k)=c_i(x_k)$ (the incumbent is always a sample point), so in particular
$\theta_k^m(x_k)=\theta(x_k)=\theta_k$; this makes each predicted decrease
$m(x_k)-m(x_k+s)$ a true model decrease from the incumbent. We denote the model
gradients and the model Jacobian by
\[
  g_k \;=\; \nabla m_f^k(x_k)\in\Rn,
  \qquad
  J_k \;=\; \bigl[\,\nabla m_{c,1}^k(x_k),\dots,\nabla m_{c,p}^k(x_k)\,\bigr]^{\!\top}\in\R^{p\times n}.
\]
These are exactly the quantities returned by \path{pounders/py/formquad.py} for a
vector-valued blackbox $F=(f,c_1,\dots,c_p)$, together with its \verb|valid| flag
(the computational witness of the fully-linear property below).

\begin{definition}[Fully-linear model, {\citealp[Def.~6.1--6.2]{CSV2009}}]
  \label{def:fl}
  A model $m$ of a function $h\in\{f,c_1,\dots,c_p\}$ is \emph{fully linear} on
  $B(x_k,\Delta_k)$ if there exist constants $\kappa_{\mathrm{ef}},
  \kappa_{\mathrm{eg}}\ge 0$, \emph{independent of $k$}, such that for all
  $y\in B(x_k,\Delta_k)$,
  \begin{align}
    \bigl|h(y)-m(y)\bigr|            &\;\le\; \kappa_{\mathrm{ef}}\,\Delta_k^2,
    \label{eq:fl-value}\\
    \bigl\|\nabla h(y)-\nabla m(y)\bigr\| &\;\le\; \kappa_{\mathrm{eg}}\,\Delta_k.
    \label{eq:fl-grad}
  \end{align}
\end{definition}

\begin{assumption}[Uniform fully-linearity]
  \label{ass:fl}
  The models $m_f^k$ and $m_{c,i}^k$ ($i=1,\dots,p$) are fully linear on
  $B(x_k,\Delta_k)$ at \emph{every} iteration, with common constants
  $\kappa_{\mathrm{ef}},\kappa_{\mathrm{eg}}$ that do not depend on $k$. Uniformity
  is delivered by $\Lambda$-poised geometry management (\citealp[Ch.~6]{CSV2009};
  \path{Algorithm_6_4.m}, \path{Model_Improvement_Step.m}), and full-linearity at
  \emph{every} iteration by running a model-improvement step whenever the incumbent
  set is not certifiably $\Lambda$-poised on $B(x_k,\Delta_k)$ (the standard DFO
  convention that lets the trust-region lemmas of Section~\ref{sec:tr-lemmas} apply
  uniformly). Consequently $\theta$ also has a fully-linear model with constants
  $\kappa_{\mathrm{ef}}^\theta,\kappa_{\mathrm{eg}}^\theta$ (Lemma~\ref{lem:L1}).
\end{assumption}

We stress the theme this note exists to make precise: \emph{building fully-linear
models of $f$ and $c$ on every iteration is necessary but not, by itself,
sufficient} for convergence. Two further ingredients are required --- a
criticality step that ties $\Delta_k$ to the criticality measures
(Section~\ref{sec:crit}), and a funnel that controls infeasibility
(Section~\ref{sec:funnel}) --- and the analysis of milestone M3 will pin down
exactly where each is used.

\section{Infeasibility measure}
\label{sec:infeas}

We measure constraint violation by the \emph{smooth} squared measure
\begin{equation}
  \label{eq:theta}
  \theta(x) \;=\; \tfrac12\bigl\|\,[\,c(x)\,]_+\,\bigr\|_2^2
           \;=\; \tfrac12\sum_{i=1}^p [\,c_i(x)\,]_+^2 ,
\end{equation}
so $\theta(x)=0 \iff x\in\feas$. We deliberately choose the squared form over the
unsquared $\|[c]_+\|$: since $t\mapsto[t]_+^2$ is $C^1$ with derivative $2[t]_+$,
\eqref{eq:theta} is \emph{continuously differentiable} (Assumption~\ref{ass:smooth})
with
\begin{equation}
  \label{eq:grad-theta}
  \nabla\theta(x) \;=\; \sum_{i=1}^p [\,c_i(x)\,]_+\,\nabla c_i(x)
                 \;=\; J(x)^{\top}[\,c(x)\,]_+ ,
\end{equation}
where $J(x)=[\nabla c_i(x)]^\top$ is the constraint Jacobian. This removes the
active/violated-boundary nonsmoothness of $\|[c]_+\|$ and makes the feasibility
analysis a \emph{smooth} trust-region minimization of $\theta$; it also matches the
measure $\tfrac12\|c\|^2$ of \citet{GouldToint2010,SampaioToint2015}, from which we
reproduce the requisite trust-region machinery (Section~\ref{sec:tr-lemmas}). Note
$\nabla\theta$ is Lipschitz on $\Lset$ (composition of Lipschitz maps), and
$\theta$ inherits fully-linear models from those of $c$ (Lemma~\ref{lem:L1}); it is
\emph{not} a penalty --- infeasibility is controlled by the funnel, never folded
into the objective.

The model counterpart of $\theta$, built from the fully-linear constraint models
$m_c^k=(m_{c,1}^k,\dots,m_{c,p}^k)$, is
\begin{equation}
  \label{eq:theta-model}
  \theta_k^m(x) \;=\; \tfrac12\bigl\|\,[\,m_c^k(x)\,]_+\,\bigr\|_2^2,
  \qquad
  \nabla\theta_k^m(x_k)=J_k^{\top}[\,m_c^k(x_k)\,]_+ ,
\end{equation}
and its linearization at $x_k$, $\ell_k(s)=\tfrac12\|[\,c(x_k)+J_k s\,]_+\|^2$, is
what the normal step reduces.

\section{Two criticality measures}
\label{sec:crit}

The method monitors two nonnegative quantities that vanish exactly at the two
kinds of stationarity we care about.

\paragraph{Feasibility criticality $\chi_k^c$.}
Because $\theta$ is now smooth, its first-order criticality is simply the model
gradient norm
\begin{equation}
  \label{eq:chi-c}
  \chi_k^c \;=\; \bigl\|\nabla\theta_k^m(x_k)\bigr\|
           \;=\; \bigl\|J_k^{\top}[\,m_c^k(x_k)\,]_+\bigr\| ,
\end{equation}
the standard measure for unconstrained minimization of $\theta_k^m$; $\chi_k^c=0$
means $x_k$ is a stationary point of the model infeasibility. (This is the model
analog of $\|\nabla\theta\|=\|J^\top[c]_+\|$, the Gould--Toint feasibility measure.)

\paragraph{Optimality / KKT criticality $\pi_k^f$.}
This measures first-order KKT stationarity for~\eqref{eq:problem} using the
\emph{model} data $(g_k,J_k)$ and the linearized active constraints. With the
model active set
$\A_k(\varepsilon)=\{\,i : c_i(x_k) \ge -\varepsilon\,\}$, write
$J_{k}^{\A}$ for the rows of $J_k$ indexed by $\A_k$. A convenient
projected-gradient form is
\begin{equation}
  \label{eq:pi-f}
  \pi_k^f \;=\;
  \Bigl|\, \min_{\substack{\|d\|\le 1,\ J_k^{\A} d\,\le\,0}}
      \ \bigl\langle g_k,\, d\bigr\rangle \,\Bigr|,
\end{equation}
the largest first-order decrease of the model objective achievable along a
direction that keeps the linearized active constraints from being violated.
Equivalently, under a constraint qualification, $\pi_k^f$ is a residual of the
KKT system
\begin{equation}
  \label{eq:kkt}
  g_k + J_k^{\top}\lambda = 0,\qquad
  \lambda\ge 0,\qquad
  \lambda_i\, c_i(x_k)=0,\quad i=1,\dots,p,
\end{equation}
i.e.\ $\pi_k^f=0$ certifies that $(x_k,\lambda)$ solves the \emph{model} KKT
conditions. The constraint qualification we invoke to pass from model KKT to true
KKT and to bound the multipliers $\lambda$ is the following.

\begin{assumption}[MFCQ at limit points]
  \label{ass:mfcq}
  At every limit point $x_*$ of the sequence of incumbents, the
  Mangasarian--Fromovitz constraint qualification holds: the gradients
  $\{\nabla c_i(x_*)\}_{i\in\A(x_*)}$ of the active constraints
  ($\A(x_*)=\{i:c_i(x_*)=0\}$) admit a direction $d$ with
  $\langle\nabla c_i(x_*),d\rangle<0$ for all $i\in\A(x_*)$.
\end{assumption}

MFCQ is equivalent to boundedness of the set of Lagrange multipliers and to
consistency of the linearized active constraints; both are used in M3 (the
normal-step decrease lemma L2 needs linearized consistency, and the passage to a
KKT limit point in T2 needs bounded multipliers).

A point $x_*$ is our target \emph{constrained (KKT) stationary point} if
$\theta(x_*)=0$ and there exists $\lambda\ge0$ with
$\nabla f(x_*)+\sum_i\lambda_i\nabla c_i(x_*)=0$ and
$\lambda_i c_i(x_*)=0$.

\section{The trust-funnel mechanism for inequalities}
\label{sec:funnel}

Rather than penalizing infeasibility, we cap it. The method maintains a
\emph{funnel bound}
\begin{equation}
  \label{eq:funnel}
  \theta_k^{\max} > 0,
  \qquad \text{maintained monotone non-increasing:}\quad
  \theta_{k+1}^{\max}\le\theta_k^{\max},
\end{equation}
and only accepts iterates that satisfy $\theta_{k+1}\le\theta_{k+1}^{\max}$. Each
iteration is classified (as in \citealp{GouldToint2010} and its DFO descendant
\citealp{SampaioToint2015}) as either

\begin{itemize}
  \item an \emph{$f$-iteration}, when the (model-predicted) objective decrease is
        the dominant achievement and feasibility is comfortably within the funnel;
        the funnel bound is left unchanged and $x_k$ moves to reduce $f$; or
  \item a \emph{$c$-iteration}, when reducing infeasibility is the priority; a
        normal step reduces the linearized violation and the funnel bound is
        tightened, driving $\theta_k\to0$.
\end{itemize}

The switch between the two is governed by comparing $\theta_k$ against a forcing
function of the optimality criticality $\pi_k^f$ (a normal step is taken only when
$\theta_k$ is large relative to $\pi_k^f$), exactly the trust-funnel switching
condition, specialized here to the inequality infeasibility
measure~\eqref{eq:theta}. This single-cap mechanism is the \emph{funnel}; a
\emph{filter} (a Pareto set of $(\theta,f)$ pairs) is an alternative acceptance
rule with the same purpose. We default to the funnel for its cleaner monotonicity
proof and note the filter as a variant.

The precise step computation (normal step $n_k$, tangential step $t_k$,
$s_k=n_k+t_k\in B(x_k,\Delta_k)$), the ratio/acceptance tests, the funnel and
$\Delta_k$ updates, and the criticality step are specified in milestone~M2; their
convergence analysis (lemmas L1--L5, theorems T1--T2) is milestone~M3.

% !TEX root = main.tex
%======================================================================
%  M1 Scoping note --- Positioning / relation to prior work + novelty
%======================================================================
\section{Positioning: why this is a genuine gap}
\label{sec:positioning}

The setting of Section~\ref{sec:problem} --- \emph{model-based} DFO with
fully-linear models of the objective \emph{and} of \emph{relaxable inequality}
constraints, aiming at KKT stationarity under MFCQ --- sits in a gap left open by
the closest lines of work. We locate each precedent precisely and identify what
it does not cover. Every claim below is stated against the source PDF in
\path{articles/}.

\paragraph{Trust-funnel lineage (derivative-based).}
\citet{GouldToint2010} introduced the trust-funnel method: equality-constrained
nonlinear programming ``without a penalty function or a filter,'' using a
normal/tangential step decomposition and a monotone infeasibility cap (the
funnel). We adopt this normal/tangential + funnel machinery
(Section~\ref{sec:funnel}) but replace derivatives with fully-linear models and
target inequalities.

\paragraph{Derivative-free trust-funnel (equalities).}
\citet{SampaioToint2015} is the primary precedent: a derivative-free trust-funnel
method that builds fully-linear models and proves first-order convergence --- but
for \emph{equality} constraints $c(x)=0$. The companion numerical study
\citep{SampaioToint2016} reports experience extending the approach to problems
``with both equality and inequality constraints and \dots\ simple bounds,''
handling inequalities by the usual device of adding slack variables and using
``neither filter nor penalty function.'' Our project differs in two ways that
matter for the theory: (i) we treat inequalities \emph{directly} through the
violation measure $\theta=\|[c]_+\|$ rather than by slack-variable
equality reformulation, so the model geometry lives in the original $x$-space that
the simulation samples; and (ii) our target is a \emph{proof} of KKT convergence
under MFCQ for that direct inequality treatment, not a numerical extension.

\paragraph{Model-based DFO with feasible inequality iterates (the contrast).}
NOWPAC \citep{AugustinMarzouk2017} and its stochastic generalization (S)NOWPAC
\citep{Menhorn2022} build fully-linear models of the objective and of
\emph{inequality} constraints inside a trust region --- superficially the closest
match. The decisive difference is that NOWPAC enforces \emph{feasible iterates}:
it adds a convex offset $h_k$ to the constraints, the ``inner boundary path,''
which keeps iterates strictly inside the feasible region and is precisely what
their convergence argument relies on \citep[\S4]{Menhorn2022}. That design
targets \emph{unrelaxable} constraints (points outside $\feas$ may not be usable)
--- the opposite of our Assumption~\ref{ass:relaxable}. Requiring feasibility
throughout forbids the infeasible model-sample points and normal-step excursions
that our funnel both permits and needs. So NOWPAC solves an adjacent but distinct
problem, and its proof technique does not transfer to the relaxable setting.

\paragraph{Feasibility-restoration alternatives.}
Inexact-restoration DFO \citep{Birgin2021} handles constraints by alternating a
feasibility-restoration phase with an optimization phase. It is a legitimate
alternative architecture for relaxable constraints and a useful contrast, but it
is not a trust-funnel method and does not build a single fully-linear model set
used jointly by a normal and a tangential step; we mention it as related rather
than as the framework we extend.

\paragraph{Constraint taxonomy and the DFO backdrop.}
\citet{LeDigabelWild2015} give the taxonomy (quantifiable/relaxable/simulation)
that names our constraint class and explains why relaxability is the enabling
property. \citet{LMW2019} survey constrained DFO and the fully-linear-model /
trust-region machinery, and \citet[Ch.~6,10,11]{CSV2009} provide the
$\Lambda$-poisedness geometry (uniform fully-linearity constants) and the
unconstrained trust-region convergence template that our criticality step and
$\kappa_{\mathrm{eg}}\Delta_k\to0$ argument specialize.

\subsection{Novelty statement}
\label{sec:novelty}

To our knowledge, no existing method provides a \emph{model-based} trust-funnel
algorithm --- with fully-linear models of \emph{both} the objective and the
\emph{inequality} constraints --- for \emph{relaxable} simulation constraints,
together with a proof of convergence to a KKT stationary point under MFCQ, and
\emph{without} a penalty parameter or a feasible-iterate requirement.

\medskip\noindent
Concretely, the contribution combines
\begin{enumerate}
  \item the DFO trust-funnel of \citet{SampaioToint2015} --- but extended from
        \emph{equalities} to \emph{inequalities} via the violation measure
        $\theta=\|[c]_+\|$ handled \emph{directly} in $x$-space;
  \item the fully-linear objective-and-inequality-constraint modeling of NOWPAC
        \citep{AugustinMarzouk2017,Menhorn2022} --- but \emph{without} its
        feasible-iterate inner-boundary requirement, using relaxability
        (Assumption~\ref{ass:relaxable}) to sample and step through infeasible
        points; and
  \item a KKT-under-MFCQ convergence guarantee, obtained without any penalty
        parameter (contrast the penalty $h$-functions in the repository, e.g.\
        \path{general_h_funs.py}, which this project deliberately does \emph{not}
        use).
\end{enumerate}

\subsection{Summary table}
\label{sec:table}

\begin{center}\small
\begin{tabular}{@{}lccccc@{}}
\toprule
Method & DF? & FL model of $c$ & Constraint type & Iterates & Convergence \\
\midrule
Gould--Toint \citeyearpar{GouldToint2010}   & no  & --- & equality            & funnel  & 1st-order \\
Sampaio--Toint \citeyearpar{SampaioToint2015} & yes & yes & equality           & funnel  & 1st-order \\
NOWPAC / (S)NOWPAC                            & yes & yes & inequality          & feasible-only & yes \\
Inexact restoration \citeyearpar{Birgin2021} & yes & (rest.) & general         & restoration & yes \\
\textbf{This project}                        & \textbf{yes} & \textbf{yes} & \textbf{relaxable ineq.} & \textbf{funnel (infeasible OK)} & \textbf{KKT/MFCQ (target)} \\
\bottomrule
\end{tabular}
\end{center}

\noindent
``FL model of $c$'' = fully-linear model of the constraints; ``DF'' =
derivative-free. The last row is the deliverable of milestones M2--M3.


% ---- M2 content ----
% !TEX root = main.tex
%======================================================================
%  M2 --- Algorithm specification
%  "DFO trust-funnel for relaxable inequalities"
%  Mirrors DFO_quad_TR/Algorithm_11_1.m (5-step) and Algorithm_10_2.m.
%======================================================================
\section{Algorithm: a DFO trust-funnel for relaxable inequalities}
\label{sec:algorithm}

This section specifies the method whose convergence is analyzed in
milestone~M3. It keeps the five-step structure of the CSV modified-DFO
algorithm as implemented in \path{DFO_quad_TR/Algorithm_11_1.m}
(criticality $\to$ step calculation $\to$ acceptance $\to$ model improvement
$\to$ radius update), and grafts onto it the normal/tangential step and funnel
machinery of \citet{GouldToint2010} and its derivative-free descendant
\citet{SampaioToint2015}, specialized to the smooth inequality infeasibility measure
$\theta=\tfrac12\|[c]_+\|^2$ of Section~\ref{sec:infeas}. The design deliberately avoids
both a penalty parameter and a feasible-iterate requirement
(Section~\ref{sec:positioning}).

\subsection{Per-iteration data and constants}
\label{sec:algo-data}

Entering iteration $k$ we hold the incumbent $x_k$, radius $\Delta_k$, funnel
bound $\theta_k^{\max}$, and the fully-linear models $m_f^k,\,m_{c,i}^k$ of
Section~\ref{sec:models} with gradient $g_k=\nabla m_f^k(x_k)$ and Jacobian
$J_k=[\nabla m_{c,i}^k(x_k)]$. Write $c_k=c(x_k)$, $f_k=f(x_k)$,
$\theta_k=\theta(x_k)=\tfrac12\|[c_k]_+\|^2$. Because $f$ and $c$ come from one vector
blackbox $F=(f,c_1,\dots,c_p)$, all $p+1$ models are produced by a single call to
\path{formquad} sharing the same $\Lambda$-poised sample set, and the
\verb|valid| flag is the computational witness of Assumption~\ref{ass:fl}.

The constants below reuse the names in \path{DFOoptions.mat} /
\path{Algorithm_11_1.m}; the funnel/step constants are new and named in the
Gould--Toint style. All are fixed over the run.
\[
\begin{array}{llll}
  \eta_0\in(0,\eta_1), & \eta_1\in(0,1) & \text{(ratio-test thresholds)} & \text{\ttfamily eta\_0, eta\_1}\\[2pt]
  \gamma\in(0,1), & \gamma_{\text{inc}}>1 & \text{(radius shrink / expand)} & \text{\ttfamily gamma, gamma\_inc}\\[2pt]
  \Delta_{\max}>0, & \Delta_{\min}>0 & \text{(radius cap / stop)} & \text{\ttfamily Delta\_max, Delta\_min}\\[2pt]
  \epsilon_c>0, & \omega\in(0,1) & \text{(criticality threshold / shrink)} & \text{\ttfamily epsilon\_c, omega}\\[2pt]
  \mu>0, & \beta>0 & \text{(criticality--radius coupling)} & \text{\ttfamily mu, beta}\\[2pt]
  r\ge1 & & \text{(poisedness-radius factor)} & \text{\ttfamily r}\\[2pt]
  \kappa_{\mathrm n}\in(0,1] & & \text{(normal-step radius fraction)} & \text{new}\\[2pt]
  \omega_{\mathrm n},\ \omega_{\mathrm t} & & \text{(normal-trigger / switching bounding functions)} & \text{new}\\[2pt]
  \kappa_\omega\in(0,1) & & \text{(trigger/switching consistency)} & \text{new}\\[2pt]
  \kappa_{\mathrm{tg}}\in(0,1) & & \text{(tangential-objective domination)} & \text{new}\\[2pt]
  \kappa_{\mathrm{nt}}\in(0,1) & & \text{(feasibility-retention relaxation)} & \text{new}\\[2pt]
  \kappa_{\mathrm{ca}}\in(0,1),\ \kappa_{\mathrm{cb}}\in(0,1) & & \text{(funnel-update constants)} & \text{new}\\[2pt]
  \kappa_{\eps}>0 & & \text{(active-set tolerance scale)} & \text{new}
\end{array}
\]
Here $\omega_{\mathrm n},\omega_{\mathrm t}:[0,\infty)\to[0,\infty)$ are continuous,
non-decreasing \emph{bounding functions} with $\omega_{\mathrm n}(t)=0\iff t=0$ and
$\omega_{\mathrm t}$ strictly increasing near $0$; the canonical choice is
$\omega_{\mathrm n}(t)=\omega_{\mathrm t}(t)=\kappa\,t^{\,q}$ for constants
$\kappa>0$, $q\ge1$.

\paragraph{Standing constant relations (used in M3).} The analysis requires the
following relations, all satisfiable by construction:
\begin{enumerate}\itemsep1pt
  \item \emph{Feasibility radius cap:} on normal-step iterations the model management
        enforces $\Delta_k\le\mu\chi_k^c$ \eqref{eq:feas-cap}, so (with L1 and
        cq-nonstat) $\chi_k^c\ge\kappa_{\mathrm m}\sqrt{2\theta_k}/(1+\mu\kappa_{\mathrm{eg}}^\theta)$:
        an over-large radius cannot make the model report a spuriously small
        feasibility criticality. (With the smooth measure this replaces the
        $\epsilon_c$-vs-$\kappa_{\mathrm m}$ relation of the unsquared draft.)
  \item \emph{Trigger/switching consistency} (\citealp[(2.31)]{GouldToint2010}):
        $\omega_{\mathrm t}(\omega_{\mathrm n}(t))\le\kappa_\omega\,t$ for all
        $t\ge0$. This couples the normal-step trigger \eqref{eq:normal-trigger} and
        the switching test \eqref{eq:switch} so that they cannot both be dormant
        while $s_k\ne0$ is productive (the residual $s_k=0$ case is the
        $\mu$-iteration of Step~3, controlled by Lemma~\ref{lem:mu-critical}).
  \item \emph{Shrinking active-set tolerance:} the model active set uses
        $\varepsilon_k=\kappa_{\eps}\,\sigma_k\to0$ (rather than a fixed
        $\varepsilon$), so $\A_k(\varepsilon_k)$ identifies the true active set at a
        limit point (needed for limiting complementarity in T2).
\end{enumerate}

\subsection{Step~1 --- Criticality step}
\label{sec:algo-crit}

Define the aggregate criticality
\begin{equation}
  \label{eq:sigma}
  \sigma_k \;=\; \max\bigl(\chi_k^c,\ \pi_k^f\bigr),
\end{equation}
with $\chi_k^c,\pi_k^f$ from \eqref{eq:chi-c}--\eqref{eq:pi-f}. This is the
constrained analog of the scalar \verb|sigma =|
$\max(\|c_{\mathrm{icb}}\|,-\lambda_{\min}(Q))$ used in
\path{Algorithm_11_1.m}: it collapses ``how nonstationary are we'' for both
feasibility and optimality into one number.

If $\sigma_k>\epsilon_c$, skip to Step~2 with the incumbent models and
$\Delta_k$. Otherwise invoke the criticality loop (the analog of
\path{Algorithm_10_2.m}): while the models are not certifiably fully linear on
$B(x_k,r\Delta_k)$ \emph{or} $r\Delta_k>\mu\,\sigma_k$, apply a
$\Lambda$-poised model-improvement step on $B(x_k,\omega^{i}\Delta_k)$,
rebuild $m_f^k,m_{c,i}^k$, recompute $\sigma_k$, and set
$\Delta_k\leftarrow\omega^{i}\Delta_k$. On exit set
\begin{equation}
  \label{eq:crit-radius}
  \Delta_k \;\leftarrow\; \min\bigl(\max(\Delta_k,\ \beta\,\sigma_k),\ \Delta_k^{\text{in}}\bigr),
\end{equation}
matching the final line of \path{Algorithm_10_2.m}. If $\Delta_k<\Delta_{\min}$,
stop (approximate stationarity). This is the step that \emph{ties $\Delta_k$ to
$\sigma_k$} and is what lets model stationarity transfer to true stationarity in
M3; it is the ingredient that ``fully-linear models on every iteration'' does not
supply on its own.

\paragraph{Feasibility radius cap on normal-step iterations.} Whenever a normal step
will be computed (Step~2a), the model-management additionally enforces
\begin{equation}
  \label{eq:feas-cap}
  \Delta_k \;\le\; \mu\,\chi_k^c ,
\end{equation}
i.e.\ the trust region is sized to the \emph{feasibility} criticality it is about to
reduce (shrinking $\Delta_k$ and rebuilding as in the criticality loop if
\eqref{eq:feas-cap} is violated). This is the feasibility analog of the standard
criticality step \citep[Ch.~10]{CSV2009}: it prevents an over-large radius from
letting the fully-linear model report a spuriously small $\chi_k^c$ while the true
infeasibility criticality is large. With \eqref{eq:feas-cap} and L1, cq-nonstat
yields $\chi_k^c\ge\kappa_{\mathrm m}\sqrt{2\theta_k}/(1+\mu\kappa_{\mathrm{eg}}^\theta)$
on every normal-step iteration (used throughout the feasibility analysis, M3).

\subsection{Step~2 --- Step calculation (normal then tangential)}
\label{sec:algo-step}

The trial step is a composite $s_k=n_k+t_k$ kept in $B(x_k,\Delta_k)$.

\paragraph{(2a) Normal step $n_k$ --- reduce linearized infeasibility.}
A normal step is computed when the incumbent infeasibility is significant
\emph{relative to the optimality measure of the previous iterate}, or when the
current point is too infeasible to pursue optimality (switching \eqref{eq:switch}
fails), i.e.\ when
\begin{equation}
  \label{eq:normal-trigger}
  \theta_k>0 \quad\text{and}\quad
  \Bigl(\ \theta_k\ge\omega_{\mathrm n}(\pi_{k-1}^f)
     \ \ \text{or}\ \ \pi_k^f\le\omega_{\mathrm t}(\theta_k)\ \Bigr),
\end{equation}
and $n_k=0$ otherwise. The second disjunct (switching failure) guarantees that
\emph{every} $c$-iteration carries a genuine normal step, which the feasibility
analysis (L8) needs; the first is the optimality-forcing trigger of
\citet[(2.30)]{GouldToint2010}, \textbf{the correct replacement for a radius-based one.}
A trigger such as $\theta_k>\kappa\Delta_k^2$ is a rate error (near a critical point
both $\theta_k\to0$ and $\Delta_k\to0$, and $\theta_k>\kappa\Delta_k^2$ can persist,
e.g.\ $\theta_k\asymp\Delta_k\gg\Delta_k^2$, never switching off); a purely
\emph{funnel}-relative trigger $\theta_k>\kappa\,\theta_k^{\max}$ can instead
\emph{stall} at an infeasible, non-$\theta$-stationary point where the funnel is
frozen. Tying the trigger to $\pi_{k-1}^f$ via $\omega_{\mathrm n}$ (with
$\omega_{\mathrm n}(t)=0\iff t=0$ and the consistency
$\omega_{\mathrm t}(\omega_{\mathrm n}(t))\le\kappa_\omega t$) forces feasibility
\emph{as optimality improves}: as $\pi_{k-1}^f\to0$, even a tiny $\theta_k$
triggers a normal step. This coupling is what drives $\theta_k\to0$ in T1. When
triggered,
$n_k$ approximately solves the convex (piecewise-linear/quadratic) trust-region
subproblem in the (smooth) linearized violation,
\begin{equation}
  \label{eq:normal-sub}
  \min_{n\in\Rn}\ \tfrac12\bigl\|\,[\,c_k+J_k n\,]_+\,\bigr\|^2
  \quad\text{s.t.}\quad \|n\|\le\kappa_{\mathrm n}\,\Delta_k,
\end{equation}
and is required to achieve a Cauchy-type decrease of the model infeasibility: there
is $\kappa_{\mathrm{nc}}>0$ with
\begin{equation}
  \label{eq:normal-cauchy}
  \theta_k - \theta_k^m(x_k+n_k)
  \;\ge\;
  \kappa_{\mathrm{nc}}\,\chi_k^c\,
  \min\!\Bigl(\tfrac{\chi_k^c}{\,1+\beta_\theta\,},\ \kappa_{\mathrm n}\Delta_k\Bigr),
\end{equation}
where $\beta_\theta$ bounds the curvature of $\theta_k^m$ (finite, since
$\|J_k\|\le\kappa_J$). This is the standard fraction-of-Cauchy decrease for the
smooth $\theta_k^m$ along the steepest-descent direction $-\nabla\theta_k^m(x_k)$
(the objective of \eqref{eq:normal-sub} is convex, so a generalized Cauchy point
attains it). The normal step is also required to be \emph{normal},
\begin{equation}
  \label{eq:normal-size}
  \|n_k\| \;\le\; \kappa_{\mathrm{nn}}\,\bigl\|[\,c_k\,]_+\bigr\|
             \;=\; \kappa_{\mathrm{nn}}\sqrt{2\theta_k}
  \qquad(\kappa_{\mathrm{nn}}>0),
\end{equation}
the standard trust-funnel requirement \citep[\S2]{GouldToint2010} that makes $n_k$
vanish as $\theta_k\to0$ (used in Lemma~\ref{lem:eventually-f}). Computationally
\eqref{eq:normal-sub} is a small bound-constrained convex QP in $(c_k,J_k)$ --- the
role played by \path{minqsw}/\path{bqmin} in the code-bridge (D5).

\paragraph{(2b) Tangential step $t_k$ --- reduce the objective model.}
Given $n_k$, the tangential step reduces the objective model along directions that
do not undo the linearized feasibility gain. Let
$\A_k=\A_k(\varepsilon_k)$ be the model active set (Section~\ref{sec:crit}) with the
\emph{shrinking} tolerance $\varepsilon_k=\kappa_{\eps}\sigma_k$, and
$N_k$ a basis for (an approximation of) the null space
$\{d:J_k^{\A}d=0\}$. We take a tangential step only when the optimality criticality
is significant relative to the current infeasibility --- the
\emph{switching condition} \citep[(2.29)]{GouldToint2010}
\begin{equation}
  \label{eq:switch}
  \pi_k^f \;>\; \omega_{\mathrm t}\bigl(\theta_k\bigr)
  \qquad\bigl(\text{equivalently } \theta_k<\omega_{\mathrm t}^{-1}(\pi_k^f)\text{: infeasibility small relative to } \pi_k^f\bigr).
\end{equation}
When \eqref{eq:switch} holds, $t_k$ approximately solves
\begin{equation}
  \label{eq:tang-sub}
  \min_{t}\ m_f^k(x_k+n_k+t)
  \quad\text{s.t.}\quad
  J_k^{\A}\,t\le 0,\quad \|n_k+t\|\le\Delta_k,\quad
  \theta_k^m\bigl(x_k+n_k+t\bigr)\le\vartheta_k,
\end{equation}
where the last constraint --- the analog of \citet[(2.28)]{GouldToint2010} for the
squared measure --- caps the model infeasibility of the composite step at
\begin{equation}
  \label{eq:retention}
  \vartheta_k \;:=\; \kappa_{\mathrm{nt}}\,\theta_k + (1-\kappa_{\mathrm{nt}})\,\theta_k^m(x_k+n_k)
  \qquad(\kappa_{\mathrm{nt}}\in(0,1)),
\end{equation}
so the tangential step cannot undo the normal step's feasibility gain by more than a
fraction. (It is feasible: $t=0$ satisfies it since $\theta_k^m(x_k+n_k)\le\vartheta_k$.)
This \emph{feasibility-retention} bound lets the composite predicted decrease inherit
the normal-step Cauchy bound:
$\pred^\theta_k=\theta_k-\theta_k^m(x_k+s_k)\ge\theta_k-\vartheta_k=(1-\kappa_{\mathrm{nt}})\bigl(\theta_k-\theta_k^m(x_k+n_k)\bigr)$,
i.e.\ $\pred^\theta_k\ge(1-\kappa_{\mathrm{nt}})\times$\eqref{eq:normal-cauchy} (used in L8).
On this step $t_k$ must deliver a fraction-of-Cauchy decrease of the objective model on the
tangent space: there is $\kappa_{\mathrm{tc}}>0$ with
\begin{equation}
  \label{eq:tang-cauchy}
  m_f^k(x_k+n_k) - m_f^k(x_k+s_k)
  \;\ge\;
  \kappa_{\mathrm{tc}}\,\pi_k^f\,
  \min\!\Bigl(\tfrac{\pi_k^f}{\,1+\|H_k\|\,},\ \Delta_k\Bigr),
\end{equation}
where $H_k$ is the objective model Hessian. If \eqref{eq:switch} fails, set
$t_k=0$ (this is a pure feasibility, or ``$c$-'', iteration). In all cases
$s_k=n_k+t_k$ with $\|s_k\|\le\Delta_k$. When $p=0$ and $\theta_k\equiv0$ this
reduces exactly to the unconstrained trust-region step of \path{Algorithm_11_1.m}
(Step~2, \verb|s_k = trust(c,Q,Delta)|).

\subsection{Step~3 --- Iteration type and acceptance}
\label{sec:algo-accept}

Define the predicted decreases
\begin{align}
  \mathrm{pred}^f_k &= m_f^k(x_k) - m_f^k(x_k+s_k), \label{eq:predf}\\
  \mathrm{pred}^\theta_k &= \theta_k - \theta_k^m(x_k+s_k), \label{eq:predtheta}
\end{align}
and the achieved decreases $\mathrm{ared}^f_k=f_k-f(x_k+s_k)$,
$\mathrm{ared}^\theta_k=\theta_k-\theta(x_k+s_k)$, computed from a single blackbox
evaluation of $F$ at the trial point $x_k+s_k$ (yielding $f$ and all $c_i$ at
once). The iteration is classified into three mutually exclusive types
(following \citealp{GouldToint2010}), the $\mu$-iteration handling the degenerate
$s_k=0$ case that a two-way $f$/$c$ split would leave undefined:

\begin{description}
\item[$\mu$-iteration] if $s_k=0$. Under the strengthened trigger
  \eqref{eq:normal-trigger} this is a \emph{model-critical} situation:
  $t_k=0$ with switching held would force $\pi_k^f=0$ (L3), contradicting
  \eqref{eq:switch}; hence switching failed ($\pi_k^f\le\omega_{\mathrm t}(\theta_k)$),
  the trigger's second disjunct fired, and $n_k=0$ then forces $\chi_k^c=0$
  (the normal Cauchy bound \eqref{eq:normal-cauchy} is null). So a $\mu$-iteration
  has $\chi_k^c=0$ and $\pi_k^f\le\omega_{\mathrm t}(\theta_k)$
  (Lemma~\ref{lem:mu-critical}); by cq-nonstat \eqref{eq:cq-nonstat}, $\chi_k^c=0$
  forces $\sqrt{2\theta_k}\le\kappa_{\mathrm{eg}}^\theta\Delta_k/\kappa_{\mathrm m}$, i.e.\
  the point is \emph{near-feasible relative to the radius}. No trial point is
  evaluated; the iteration is routed through the criticality step (Step~1), the funnel
  is unchanged, and $x_{k+1}=x_k$. In particular, a $\mu$-iteration \emph{cannot occur
  at a point with $\theta_k\ge\tau>0$ and
  $\Delta_k<\kappa_{\mathrm m}\sqrt{2\tau}/\kappa_{\mathrm{eg}}^\theta$} --- the fact the
  feasibility analysis (L8) relies on.
\item[$f$-iteration] if $s_k\ne0$, a tangential step was computed
  (\eqref{eq:switch} held), and the full objective decrease retains a fraction of the
  \emph{tangential} objective decrease,
  \begin{equation}
    \label{eq:f-domination}
    \mathrm{pred}^f_k \;\ge\; \kappa_{\mathrm{tg}}\,\delta^{f,t}_k,
    \qquad \delta^{f,t}_k := m_f^k(x_k+n_k)-m_f^k(x_k+s_k)
  \end{equation}
  (the \citet[(2.33)]{GouldToint2010} domination test; $\delta^{f,t}_k$ is the
  quantity bounded below by L3, so \eqref{eq:f-domination} transfers that bound to
  $\mathrm{pred}^f_k$ --- this is what the feasibility path-length lemma L6 uses,
  and it replaces the earlier, incorrect comparison of $\mathrm{pred}^f_k$ to
  $\mathrm{pred}^\theta_k$). Acceptance uses the objective ratio
  \begin{equation}
    \label{eq:rhof}
    \rho^f_k \;=\; \frac{\mathrm{ared}^f_k}{\mathrm{pred}^f_k}.
  \end{equation}
  The trial point is accepted iff it respects the funnel and the ratio passes:
  \begin{equation}
    \label{eq:accept-f}
    \theta(x_k+s_k)\le\theta_k^{\max}
    \quad\text{and}\quad
    \bigl(\rho^f_k\ge\eta_1,\ \text{or}\ (\rho^f_k>\eta_0\ \text{and models valid})\bigr).
  \end{equation}
  The funnel bound is left unchanged: $\theta_{k+1}^{\max}=\theta_k^{\max}$.
\item[$c$-iteration] otherwise ($s_k\ne0$ and \eqref{eq:f-domination} fails or no
  tangential step was computed). \emph{A $c$-iteration always carries a genuine
  normal step} ($n_k\ne0$ whenever $\chi_k^c>0$): if switching \eqref{eq:switch}
  failed, the trigger \eqref{eq:normal-trigger} computed $n_k$ directly; if switching
  held but \eqref{eq:f-domination} failed, then $\mathrm{pred}^f_k<\kappa_{\mathrm{tg}}\delta^{f,t}_k\le\delta^{f,t}_k$,
  which (as $\mathrm{pred}^f_k=\delta^{f,t}_k+[m_f^k(x_k)-m_f^k(x_k+n_k)]$) forces
  $m_f^k(x_k)-m_f^k(x_k+n_k)<0$, i.e.\ $n_k\ne0$. This closes the gap that a pure
  tangential step could be classified $c$ without a normal step. Acceptance uses the
  feasibility ratio
  \begin{equation}
    \label{eq:rhoc}
    \rho^c_k \;=\; \frac{\mathrm{ared}^\theta_k}{\mathrm{pred}^\theta_k},
  \end{equation}
  well defined since the normal step gives, via retention \eqref{eq:retention},
  $\mathrm{pred}^\theta_k\ge(1-\kappa_{\mathrm{nt}})\times\eqref{eq:normal-cauchy}>0$.
  The trial point is accepted iff $\rho^c_k\ge\eta_1$. On a successful $c$-iteration
  the funnel is tightened (Step~4b).
\end{description}

On acceptance, $x_{k+1}=x_k+s_k$; on rejection, $x_{k+1}=x_k$. As in
\path{Algorithm_11_1.m}, the trial point (and its already-paid blackbox value) is
\emph{always} added to the sample set regardless of acceptance, since evaluations
are expensive.

\subsection{Step~4 --- Model improvement and funnel update}
\label{sec:algo-improve}

\paragraph{(4a) Model improvement.} If the iteration was unsuccessful
($\rho^{\bullet}_k<\eta_1$) or the step was negligible
($\|s_k\|/\Delta_k$ below \path{tol\_step\_ratio}), run a $\Lambda$-poised
model-improvement step on $B(x_k,r\Delta_k)$ and rebuild the models --- exactly
Step~4 of \path{Algorithm_11_1.m} via \path{Model_Improvement_Step.m}.

\paragraph{(4b) Funnel update.} On a successful $c$-iteration, tighten the funnel
by the Gould--Toint rule
\begin{equation}
  \label{eq:funnel-update}
  \theta_{k+1}^{\max}
  = \max\Bigl(\kappa_{\mathrm{ca}}\,\theta_k^{\max},\
                  \theta_{k+1} + \kappa_{\mathrm{cb}}\,(\theta_k^{\max}-\theta_{k+1})\Bigr),
\end{equation}
which keeps $\theta_{k+1}<\theta_{k+1}^{\max}\le\theta_k^{\max}$ (strict room
below the cap, monotone non-increasing cap). On $f$-iterations and unsuccessful
iterations, $\theta_{k+1}^{\max}=\theta_k^{\max}$. This monotonicity is lemma~L4
of M3 and is the mechanism that forces $\theta_k\to0$.

\subsection{Step~5 --- Trust-region radius update}
\label{sec:algo-radius}

Mirroring Step~5 of \path{Algorithm_11_1.m}:
\begin{equation}
  \label{eq:radius-update}
  \Delta_{k+1} =
  \begin{cases}
    \min(\gamma_{\text{inc}}\Delta_k,\ \Delta_{\max})
      & \text{very successful } (\rho^{\bullet}_k\ge\eta_1) \text{ and } \|s_k\|\approx\Delta_k,\\[4pt]
    \gamma\,\Delta_k
      & \text{unsuccessful } (\rho^{\bullet}_k<\eta_1) \text{ and models valid on } B(x_k,r\Delta_k),\\[4pt]
    \Delta_k & \text{otherwise,}
  \end{cases}
\end{equation}
where $\rho^{\bullet}_k$ is $\rho^f_k$ on $f$-iterations and $\rho^c_k$ on
$c$-iterations. Increment $k$ and return to Step~1.

\subsection{Consolidated pseudocode}
\label{sec:algo-pseudocode}

\begin{quote}\small
\noindent\textbf{Algorithm~1 (DFO trust-funnel for relaxable inequalities).}\\
\textit{Given} feasible $x_0$ (Assumption~\ref{ass:feasstart}), $\Delta_0$,
$\theta_0^{\max}\ge\theta_0$, constants of Section~\ref{sec:algo-data}.\\
\textbf{For} $k=0,1,2,\dots$:
\begin{enumerate}\itemsep2pt
  \item \textbf{Models.} Build fully-linear $m_f^k,\{m_{c,i}^k\}$ on
        $B(x_k,\Delta_k)$; form $g_k,J_k$; compute $\chi_k^c,\pi_k^f$,
        $\sigma_k=\max(\chi_k^c,\pi_k^f)$.
  \item \textbf{Criticality.} If $\sigma_k\le\epsilon_c$: run the
        \path{Algorithm_10_2} loop (improve poisedness, shrink $\Delta_k$ toward
        $\mu\sigma_k$, apply \eqref{eq:crit-radius}); if $\Delta_k<\Delta_{\min}$
        \textbf{stop}.
  \item \textbf{Normal step.} If $\theta_k\ge\omega_{\mathrm n}(\pi_{k-1}^f)$,
        compute $n_k$ solving \eqref{eq:normal-sub} with decrease
        \eqref{eq:normal-cauchy}; else $n_k=0$.
  \item \textbf{Tangential step.} If switching \eqref{eq:switch} holds, compute
        $t_k$ solving \eqref{eq:tang-sub} with decrease \eqref{eq:tang-cauchy};
        else $t_k=0$. Set $s_k=n_k+t_k$ ($\|s_k\|\le\Delta_k$).
  \item \textbf{Classify \& evaluate.} If $s_k=0$: \textbf{$\mu$-iteration} ---
        criticality/geometry refresh, $x_{k+1}=x_k$, funnel unchanged, go to Step~7.
        Else evaluate $F(x_k+s_k)$ and classify $f$- vs.\ $c$-iteration (Step~3);
        accept/reject by \eqref{eq:accept-f}/\eqref{eq:rhoc}.
  \item \textbf{Improve \& funnel.} Model improvement (4a) if needed; funnel
        update \eqref{eq:funnel-update} on a successful $c$-iteration.
  \item \textbf{Radius.} Update $\Delta_{k+1}$ by \eqref{eq:radius-update}.
\end{enumerate}
\end{quote}

\paragraph{Design invariants (to be discharged in M3).}
(i) $\theta_k\le\theta_k^{\max}$ for all $k$ and $\{\theta_k^{\max}\}$ is
non-increasing; (ii) every accepted step lies in $B(x_k,\Delta_k)$; (iii) with
$p=0$ (so $\theta_k\equiv0$, $n_k\equiv0$, always $f$-iterations) Algorithm~1
collapses to \path{Algorithm_11_1.m}; (iv) rewriting an equality
$c^{\mathrm{eq}}(x)=0$ as $\{c^{\mathrm{eq}}\le0,\ -c^{\mathrm{eq}}\le0\}$
recovers the equality trust-funnel of \citet{SampaioToint2015}. These invariants
are the sanity checks the convergence analysis (M3) must respect.


% ---- M3 content ----
% !TEX root = main.tex
%======================================================================
%  M3 --- Convergence analysis
%  Lemmas L1--L5, Theorems T1 (feasibility) and T2 (KKT under MFCQ),
%  plus limiting-case validations.
%======================================================================
\section{Convergence analysis}
\label{sec:convergence}

Throughout, Assumptions~\ref{ass:relaxable}--\ref{ass:mfcq} hold. We write
$\{x_k\}$ for the incumbents produced by Algorithm~1, $\S$ for the index set of
\emph{successful} iterations (those on which $x_{k+1}=x_k+s_k$), and we say a model
quantity is ``valid at $k$'' when the models are fully linear on $B(x_k,\Delta_k)$
(Assumption~\ref{ass:fl}). Two consequences of the standing assumptions are used
repeatedly and recorded first.

\begin{proposition}[Boundedness]
  \label{prop:bounded}
  $f$ is bounded below on the working region $\Lset$
  (Assumption~\ref{ass:levelset}), and there is $\kappa_H\ge0$ with
  $\|H_k\|\le\kappa_H$ and $\|J_k\|\le\kappa_J$ for all $k$ (model Hessian of $f$
  and model Jacobian of $c$ are uniformly bounded).
\end{proposition}
\begin{proof}
  $\Lset$ is bounded (Assumption~\ref{ass:levelset}) and $f,c\in C^1$
  (Assumption~\ref{ass:smooth}), so $f$ attains a finite minimum on $\overline{\Lset}$.
  Fully-linear gradients satisfy $\|g_k\|\le\|\nabla f(x_k)\|+\kappa_{\mathrm{eg}}\Delta_k$
  and likewise for $J_k$; since $\nabla f,\nabla c$ are bounded on $\Lset$ and
  $\Delta_k\le\Delta_{\max}$, both are uniformly bounded. Minimum-Frobenius-norm
  model Hessians (as returned by \path{formquad}) are bounded by the same data.
\end{proof}

\noindent\emph{Standing constant relations.} The analysis uses the relations fixed
in Section~\ref{sec:algo-data}: the active-set tolerance shrinks,
$\varepsilon_k=\kappa_{\eps}\sigma_k\to0$ (for limiting complementarity in T2); and,
for the smooth measure, Lemma~\ref{lem:cq-nonstat} supplies a feasibility-criticality
lower bound $\chi_k^c\ge\kappa_{\mathrm m}\sqrt{2\theta_k}-\kappa_{\mathrm{eg}}^\theta\Delta_k$
whose ``in particular'' clause is invoked with $\varrho=\tau/2$ (a fixed positive
constant in the T1 contradiction). No lower bound on $\epsilon_c$ relative to the
(unknown) infeasibility level is needed: the radius controls come from the reproduced
lemmas of Section~\ref{sec:tr-lemmas} and the criticality step (L5).

\subsection{Supporting lemmas}
\label{sec:lemmas}

\paragraph{L1 --- Model-error propagation (objective, constraints, and $\theta$).}

\begin{lemma}[L1]
  \label{lem:L1}
  Let the models be valid at $k$. Then for all $y\in B(x_k,\Delta_k)$,
  \begin{align}
    |f(y)-m_f^k(y)|            &\le \kappa_{\mathrm{ef}}\Delta_k^2, &
    \|\nabla f(y)-\nabla m_f^k(y)\| &\le \kappa_{\mathrm{eg}}\Delta_k, \label{eq:L1f}\\
    |c_i(y)-m_{c,i}^k(y)|      &\le \kappa_{\mathrm{ef}}\Delta_k^2, &
    \|\nabla c_i(y)-\nabla m_{c,i}^k(y)\| &\le \kappa_{\mathrm{eg}}\Delta_k, \label{eq:L1c}
  \end{align}
  and, with the smooth $\theta=\tfrac12\|[c]_+\|^2$ and its model
  $\theta_k^m=\tfrac12\|[m_c^k]_+\|^2$, $\theta$ has a fully-linear model on
  $B(x_k,\Delta_k)$:
  \begin{equation}
    \label{eq:L1theta}
    |\theta(y)-\theta_k^m(y)| \le \kappa_{\mathrm{ef}}^{\theta}\Delta_k^2,
    \qquad
    \|\nabla\theta(y)-\nabla\theta_k^m(y)\| \le \kappa_{\mathrm{eg}}^{\theta}\Delta_k,
  \end{equation}
  with $\kappa_{\mathrm{ef}}^{\theta}=\sqrt p\,C_\theta\kappa_{\mathrm{ef}}$ and
  $\kappa_{\mathrm{eg}}^{\theta}=\sqrt p\,(C_\theta\kappa_{\mathrm{eg}}+\kappa_J\kappa_{\mathrm{ef}}\Delta_{\max})$,
  where $C_\theta:=\sup_{\Lset}\|[c]_+\|+\sqrt p\,\kappa_{\mathrm{ef}}\Delta_{\max}^2<\infty$.
\end{lemma}
\begin{proof}
  \eqref{eq:L1f}--\eqref{eq:L1c} are Definition~\ref{def:fl} applied to $f$ and each
  $c_i$. For \eqref{eq:L1theta}, write $u=[c(y)]_+$, $v=[m_c^k(y)]_+$; the map
  $z\mapsto[z]_+$ is $1$-Lipschitz, so $\|u-v\|\le\|c(y)-m_c^k(y)\|\le\sqrt p\,\kappa_{\mathrm{ef}}\Delta_k^2$
  by \eqref{eq:L1c}. Then
  $|\theta-\theta_k^m|=\tfrac12|\,\|u\|^2-\|v\|^2\,|=\tfrac12|\langle u+v,u-v\rangle|\le\tfrac12(\|u\|+\|v\|)\|u-v\|\le C_\theta\|u-v\|$,
  giving the value bound. For the gradient, $\nabla\theta(y)=J(y)^\top[c(y)]_+$ and
  $\nabla\theta_k^m(y)=J_k^\top[m_c^k(y)]_+$; adding and subtracting $J_k^\top[c(y)]_+$
  and using \eqref{eq:L1c}, $\|[c(y)]_+\|\le C_\theta$, and $\|J_k\|\le\kappa_J$ gives
  the $O(\Delta_k)$ bound. All constants are $k$-independent (Assumption~\ref{ass:fl},
  Proposition~\ref{prop:bounded}).
\end{proof}

\noindent The point of L1 is that \emph{the smooth infeasibility measure inherits a
fully-linear model} --- $O(\Delta_k^2)$ in value and $O(\Delta_k)$ in gradient ---
from the constraint models. This lets the feasibility analysis treat $\theta$ as a
smooth merit with a fully-linear model, so the reproduced trust-region lemmas of
Section~\ref{sec:tr-lemmas} apply to it verbatim; it is why fully-linear models of
the \emph{constraints} (not merely the objective) are required.

\paragraph{L2 --- Normal-step decrease (feasibility).}

\begin{lemma}[L2]
  \label{lem:L2}
  Suppose a normal step is computed at $k$ and $n_k$ is a generalized Cauchy point of
  the convex subproblem \eqref{eq:normal-sub}. Then, with $\beta_\theta$ a curvature
  bound for $\theta_k^m$,
  \begin{equation}
    \label{eq:L2}
    \theta_k - \theta_k^m(x_k+n_k)
      \;\ge\; \kappa_{\mathrm{nc}}\,\chi_k^c\,
              \min\!\Bigl(\tfrac{\chi_k^c}{1+\beta_\theta},\ \kappa_{\mathrm n}\Delta_k\Bigr),
    \qquad \kappa_{\mathrm{nc}}=\tfrac12 .
  \end{equation}
\end{lemma}
\begin{proof}
  This is the classical fraction-of-Cauchy decrease for the \emph{smooth} convex
  $\theta_k^m$ minimized along its steepest-descent direction $-\nabla\theta_k^m(x_k)$
  inside $\|n\|\le\kappa_{\mathrm n}\Delta_k$: with criticality
  $\chi_k^c=\|\nabla\theta_k^m(x_k)\|$ and curvature $\le\beta_\theta$, the Cauchy
  point gives $\theta_k^m(x_k)-\theta_k^m(x_k+n_k^{\mathrm C})\ge\tfrac12\chi_k^c\min(\chi_k^c/\beta_\theta,\kappa_{\mathrm n}\Delta_k)$
  (\citealp[Ch.~6.3]{ConnGouldToint2000}; the same bound underlies (TR-acc) of
  Section~\ref{sec:tr-lemmas}). A generalized Cauchy point does at least this well.
  $\beta_\theta$ is finite: $\theta_k^m=\tfrac12\|[m_c^k]_+\|^2$ has Hessian (where
  differentiable) $J_k^\top D J_k+\sum_i[m_{c,i}^k]_+\nabla^2 m_{c,i}^k$ bounded by
  Proposition~\ref{prop:bounded}.
\end{proof}

\noindent \emph{Role of MFCQ.} \eqref{eq:L2} is vacuous unless $\chi_k^c>0$, i.e.\
unless $\theta_k^m$ has nonzero gradient. MFCQ (Assumption~\ref{ass:mfcq}) is exactly
what guarantees $\|\nabla\theta\|>0$ at infeasible points; this is made precise for the
smooth measure in Lemma~\ref{lem:cq-nonstat} and is the pivot of the feasibility
theorem~T1.

\paragraph{L3 --- Tangential-step decrease (optimality).}

\begin{lemma}[L3]
  \label{lem:L3}
  Suppose the switching condition \eqref{eq:switch} holds at $k$ and $t_k$ is a
  (projected) generalized Cauchy point of \eqref{eq:tang-sub}. Then
  \begin{equation}
    \label{eq:L3}
    m_f^k(x_k+n_k) - m_f^k(x_k+s_k)
      \;\ge\; \kappa_{\mathrm{tc}}\,\pi_k^f\,
              \min\!\Bigl(\tfrac{\pi_k^f}{1+\kappa_H},\ \Delta_k\Bigr),
    \qquad \kappa_{\mathrm{tc}}=\tfrac12 .
  \end{equation}
\end{lemma}
\begin{proof}
  This is the fraction-of-Cauchy-decrease bound for a quadratic model minimized
  over the polyhedral cone $\{t:J_k^{\A}t\le0\}$ intersected with the trust region,
  with $\pi_k^f$ the projected-gradient criticality \eqref{eq:pi-f}. The argument is
  the constrained Cauchy-point estimate of \citet[Ch.~10]{CSV2009} (their
  unconstrained $\|g_k\|$ replaced by the projected measure $\pi_k^f$ and the line
  search restricted to the feasible cone), using $\|H_k\|\le\kappa_H$ from
  Proposition~\ref{prop:bounded}. When $\A_k=\varnothing$ the cone is all of $\Rn$,
  $\pi_k^f=\|g_k\|$, and \eqref{eq:L3} is precisely the unconstrained Cauchy decrease
  of \path{Algorithm_11_1.m}.
\end{proof}

\paragraph{L4 --- Funnel monotonicity.}

\begin{lemma}[L4]
  \label{lem:L4}
  The funnel sequence $\{\theta_k^{\max}\}$ is non-increasing and bounded below by
  $0$, hence convergent to some $\theta_*^{\max}\ge0$. Every iterate satisfies
  $\theta_k\le\theta_k^{\max}$, and on every successful $c$-iteration
  \begin{equation}
    \label{eq:L4-gap}
    \theta_k^{\max}-\theta_{k+1}^{\max}
      \;\ge\; (1-\kappa_{\mathrm{cb}})\,(\theta_k^{\max}-\theta_{k+1}).
  \end{equation}
\end{lemma}
\begin{proof}
  By construction \eqref{eq:funnel-update} fires only on successful $c$-iterations
  and returns $\theta_{k+1}^{\max}=\max(\kappa_{\mathrm{ca}}\theta_k^{\max},\,
  \theta_{k+1}+\kappa_{\mathrm{cb}}(\theta_k^{\max}-\theta_{k+1}))$; on all other
  iterations $\theta_{k+1}^{\max}=\theta_k^{\max}$. Since a successful $c$-iteration
  has $\rho^c_k\ge\eta_1>0$, hence $\ared^\theta_k>0$ and $\theta_{k+1}<\theta_k
  \le\theta_k^{\max}$, both branches of the max are $\le\theta_k^{\max}$
  (as $\kappa_{\mathrm{ca}},\kappa_{\mathrm{cb}}\in(0,1)$), so the sequence is
  non-increasing; it is bounded below by $0$, so it converges. Acceptance rules
  \eqref{eq:accept-f}/\eqref{eq:rhoc} never admit a point with $\theta>\theta_k^{\max}$
  (the $f$-branch enforces it explicitly; the $c$-branch decreases $\theta$ from
  $\theta_k\le\theta_k^{\max}$), giving $\theta_k\le\theta_k^{\max}$ for all $k$.
  Finally, from the second argument of the max,
  $\theta_{k+1}^{\max}\le\theta_{k+1}+\kappa_{\mathrm{cb}}(\theta_k^{\max}-\theta_{k+1})$,
  which rearranges to \eqref{eq:L4-gap}.
\end{proof}

\subsection{Reproduced single-region trust-region lemmas}
\label{sec:tr-lemmas}

We reproduce, self-contained, the standard single-region trust-region results
(\citealp[Ch.~6.4]{ConnGouldToint2000}; \citealp[Ch.~10]{CSV2009}) in the exact
abstract form the constrained analysis invokes, so nothing here is imported. The
\emph{abstract setting} is one iteration acting on a merit $\varphi$ (later $f$ or
$\theta$) via a model $m$ fully linear on $B(x_k,\Delta_k)$ with value error
$\le\kappa_{\mathrm{ef}}^\varphi\Delta_k^2$ (Definition~\ref{def:fl}); a step
$\|s_k\|\le\Delta_k$ with model decrease obeying a \emph{fraction-of-Cauchy} bound
\begin{equation}
  \label{eq:fcd}
  \mathrm{pred}_k := m(x_k)-m(x_k+s_k)\;\ge\;\kappa_{\mathrm{fcd}}\,\xi_k\,
    \min\!\Bigl(\tfrac{\xi_k}{1+\kappa_H^\varphi},\,\Delta_k\Bigr)
\end{equation}
for a criticality measure $\xi_k\ge0$ and curvature bound $\kappa_H^\varphi$; ratio
$\rho_k=\mathrm{ared}_k/\mathrm{pred}_k$ with $\mathrm{ared}_k=\varphi(x_k)-\varphi(x_k+s_k)$;
acceptance iff $\rho_k\ge\eta_1$; and the radius update \eqref{eq:radius-update}
(shrink by $\gamma$ only on unsuccessful iterations; never shrink on a successful
one). All of L2, L3 supply \eqref{eq:fcd} (for $\theta$ with $\xi=\chi^c$, and for
$f$ with $\xi=\pi^f$ respectively).

\begin{lemma}[TR-acc --- accuracy implies success]
  \label{lem:tr-acc}
  In the abstract setting, if
  $\Delta_k\le\min\!\bigl(\tfrac{\xi_k}{1+\kappa_H^\varphi},\,
     \tfrac{\kappa_{\mathrm{fcd}}(1-\eta_1)\xi_k}{2\kappa_{\mathrm{ef}}^\varphi}\bigr)$
  then $\rho_k\ge\eta_1$ (the iteration is successful).
\end{lemma}
\begin{proof}
  Since $x_k,x_k+s_k\in B(x_k,\Delta_k)$, full-linearity gives
  $|\mathrm{ared}_k-\mathrm{pred}_k|\le 2\kappa_{\mathrm{ef}}^\varphi\Delta_k^2$. The
  first bound on $\Delta_k$ makes the $\min$ in \eqref{eq:fcd} equal $\Delta_k$, so
  $\mathrm{pred}_k\ge\kappa_{\mathrm{fcd}}\xi_k\Delta_k>0$. Hence
  $|\rho_k-1|\le 2\kappa_{\mathrm{ef}}^\varphi\Delta_k^2/(\kappa_{\mathrm{fcd}}\xi_k\Delta_k)
   =2\kappa_{\mathrm{ef}}^\varphi\Delta_k/(\kappa_{\mathrm{fcd}}\xi_k)\le 1-\eta_1$
  by the second bound, so $\rho_k\ge\eta_1$.
\end{proof}

\begin{lemma}[TR1 --- radius lower bound at non-critical points]
  \label{lem:tr1}
  If $\xi_j\ge\varsigma>0$ for every iteration $j$ in an index interval $[k_0,k_1]$
  on which the model is fully linear, then $\Delta_j\ge\Delta_{\mathrm{lb}}(\varsigma):=
  \gamma\min\!\bigl(\tfrac{\varsigma}{1+\kappa_H^\varphi},
    \tfrac{\kappa_{\mathrm{fcd}}(1-\eta_1)\varsigma}{2\kappa_{\mathrm{ef}}^\varphi},\Delta_{k_0}\bigr)$
  for all $j\in[k_0,k_1]$.
\end{lemma}
\begin{proof}
  Let $\bar\Delta(\varsigma)=\min(\tfrac{\varsigma}{1+\kappa_H^\varphi},
  \tfrac{\kappa_{\mathrm{fcd}}(1-\eta_1)\varsigma}{2\kappa_{\mathrm{ef}}^\varphi})$. If
  $\Delta_j\le\bar\Delta(\varsigma)\le\bar\Delta(\xi_j)$ then TR-acc makes iteration
  $j$ successful, so $\Delta_{j+1}\ge\Delta_j$ (\eqref{eq:radius-update} never shrinks
  on success). Thus the radius can drop below $\bar\Delta(\varsigma)$ only by a single
  $\gamma$-shrink from a value $>\bar\Delta(\varsigma)$, landing $\ge\gamma\bar\Delta(\varsigma)$;
  and it starts at $\Delta_{k_0}$. Hence $\Delta_j\ge\gamma\min(\bar\Delta(\varsigma),\Delta_{k_0})$
  throughout.
\end{proof}

\begin{lemma}[TR2 --- liminf criticality]
  \label{lem:tr2}
  Suppose $\varphi$ is bounded below on $\Lset$ and is non-increasing over the
  iterations under consideration (each successful one decreasing it by $\ge\eta_1\mathrm{pred}$,
  the others leaving $x$ fixed). Then $\liminf_k\xi_k=0$.
\end{lemma}
\begin{proof}
  Suppose $\xi_k\ge\varsigma>0$ for all large $k$. If there are infinitely many
  successful iterations, each decreases $\varphi$ by
  $\ge\eta_1\kappa_{\mathrm{fcd}}\varsigma\min(\varsigma/(1+\kappa_H^\varphi),\Delta_k)
  \ge\eta_1\kappa_{\mathrm{fcd}}\varsigma\min(\varsigma/(1+\kappa_H^\varphi),\Delta_{\mathrm{lb}}(\varsigma))>0$
  (TR1), so $\varphi\to-\infty$, contradicting boundedness. If instead there are only
  finitely many successful iterations, then past the last one every iteration is
  unsuccessful, so \eqref{eq:radius-update} shrinks $\Delta_k\to0$; but once
  $\Delta_k\le\bar\Delta(\varsigma)$, TR-acc forces success --- a contradiction. Either
  way $\xi_k\ge\varsigma$ eventually is impossible, so $\liminf_k\xi_k=0$.
\end{proof}

\begin{lemma}[TR-lim --- full criticality limit under the criticality step]
  \label{lem:tr-lim}
  If, in addition, the criticality step ties $\Delta_k$ to $\xi_k$
  (Lemma~\ref{lem:L5}: $\Delta_k\ge\kappa_\Delta\xi_k$ near non-critical points and
  $\Delta_k\to0$ only as $\xi_k\to0$), then $\lim_{k\to\infty}\xi_k=0$.
\end{lemma}
\begin{proof}
  Standard ``eventually-successful'' argument of \citet[Thm~6.4.6]{ConnGouldToint2000}
  / \citet[Ch.~10]{CSV2009}: were $\xi_k\ge2\varsigma$ along a subsequence while
  (by TR2) $\xi_k<\varsigma$ infinitely often, the criticality change between a
  $\xi<\varsigma$ index and the next $\xi\ge2\varsigma$ index is bounded below, forcing
  a total $\varphi$-decrease that (via TR1 and $\varphi$ bounded below) can occur only
  finitely often --- a contradiction. Hence $\lim\xi_k=0$.
\end{proof}

\paragraph{L5 --- Finiteness of the criticality step.}

\begin{lemma}[L5]
  \label{lem:L5}
  If $\sigma_k>0$, the criticality loop of Step~2 (Section~\ref{sec:algo-crit})
  terminates after finitely many model-improvement steps, and on exit
  \begin{equation}
    \label{eq:L5}
    \Delta_k \;\ge\; \min\bigl(\Delta_k^{\text{in}},\ \beta\,\sigma_k\bigr)
    \qquad\text{and}\qquad
    \Delta_k \;\le\; \mu\,\sigma_k \ \text{ or } \ \Delta_k=\Delta_k^{\text{in}} .
  \end{equation}
  Consequently there is $\kappa_\Delta>0$ with $\Delta_k\ge\kappa_\Delta\,\sigma_k$
  whenever the criticality step is entered.
\end{lemma}
\begin{proof}
  For a $\Lambda$-poised set, finitely many model-improvement steps produce a
  fully-linear model on any given radius \citep[Ch.~6, Theorem on model
  improvement]{CSV2009}. The loop shrinks the radius geometrically by $\omega\in(0,1)$
  and exits once the model is fully linear on $B(x_k,r\Delta_k)$ \emph{and}
  $r\Delta_k\le\mu\sigma_k$; since $\omega^i\Delta_k^{\text{in}}\to0<\mu\sigma_k/r$
  (as $\sigma_k>0$) and each radius needs only finitely many improvement steps, the
  loop terminates. The exit assignment \eqref{eq:crit-radius},
  $\Delta_k=\min(\max(\Delta_k,\beta\sigma_k),\Delta_k^{\text{in}})$, gives the lower
  bound in \eqref{eq:L5}; the exit test gives the upper alternative. For the final
  bound, either $\Delta_k=\Delta_k^{\text{in}}$ (the loop did not shrink, so
  $\Delta_k\ge\min(\Delta_k^{\text{in}},\beta\sigma_k)$ trivially) or the exit
  assignment gave $\Delta_k\ge\beta\sigma_k$; in both cases
  $\Delta_k\ge\beta\sigma_k$ whenever the criticality step is entered
  ($\sigma_k\le\epsilon_c$, so $\beta\sigma_k\le\beta\epsilon_c$ is the binding
  branch), i.e.\ $\kappa_\Delta=\beta$ works. This is the inequality-funnel analog of the criticality-step lemma of
  \citet[Ch.~10]{CSV2009} realized by \path{Algorithm_10_2.m}.
\end{proof}

\noindent Lemma~\ref{lem:L5} is the ingredient promised in
Section~\ref{sec:algo-crit}: it \emph{ties $\Delta_k$ to the criticality
$\sigma_k$}. Together with L1 ($\|\nabla f-g_k\|\le\kappa_{\mathrm{eg}}\Delta_k$),
it is what transfers model stationarity to true stationarity --- the step that
fully-linear models alone do not provide.

\paragraph{$\mu$-iteration criticality.}

\begin{lemma}[$\mu$-iterations are model-feasibility-critical]
  \label{lem:mu-critical}
  Under the strengthened trigger \eqref{eq:normal-trigger}, a $\mu$-iteration
  ($s_k=0$) has
  \[
    \chi_k^c=0 \qquad\text{and}\qquad \pi_k^f\le\omega_{\mathrm t}(\theta_k).
  \]
  Consequently (Lemma~\ref{lem:cq-nonstat}, contrapositive) a $\mu$-iteration occurs
  only at points that are \emph{near-feasible relative to the radius},
  $\theta_k\le(\kappa_{\mathrm{eg}}^\theta\Delta_k/\kappa_{\mathrm m})^2/2$; in
  particular no $\mu$-iteration has $\theta_k\ge\tau>0$ with
  $\Delta_k<\kappa_{\mathrm m}\sqrt{2\tau}/\kappa_{\mathrm{eg}}^\theta$.
\end{lemma}
\begin{proof}
  $s_k=0$ forces $t_k=0$ and $n_k=0$. If switching \eqref{eq:switch} held, the
  tangential Cauchy decrease \eqref{eq:tang-cauchy} is null only when $\pi_k^f=0$,
  contradicting $\pi_k^f>\omega_{\mathrm t}(\theta_k)\ge0$; hence switching failed,
  $\pi_k^f\le\omega_{\mathrm t}(\theta_k)$. Then the trigger's second disjunct fired,
  so a normal step was computed; $n_k=0$ makes the normal Cauchy bound
  \eqref{eq:normal-cauchy} null, forcing $\chi_k^c=0$. By
  Lemma~\ref{lem:cq-nonstat}, $0=\chi_k^c\ge\kappa_{\mathrm m}\sqrt{2\theta_k}-\kappa_{\mathrm{eg}}^\theta\Delta_k$,
  i.e.\ $\sqrt{2\theta_k}\le\kappa_{\mathrm{eg}}^\theta\Delta_k/\kappa_{\mathrm m}$, which
  is the stated near-feasibility bound.
\end{proof}

\noindent The lemma is what lets L8 exclude $\mu$-iterations (which shrink $\Delta$)
at infeasible points: a $\mu$-iteration is model-feasibility-critical, and under
extended MFCQ that means near-feasible relative to the radius. Its residual
infeasibility, if any, is $O(\Delta_k^2)$ and handled by the funnel drive of T1.

\subsection{Feasibility of limit points (Theorem T1)}
\label{sec:T1}

The feasibility argument rests on the observation that MFCQ makes an infeasible
point a \emph{non}-stationary point of the infeasibility measure $\theta$. We
state the precise form used. For $x$ with $\theta(x)>0$ let the
\emph{positively-active set} be $\A^+(x)=\{i:c_i(x)\ge0\}$ (constraints that are
violated or exactly active); these are the ones contributing to $\theta$.

\begin{assumption}[Extended MFCQ on the working region]
  \label{ass:emfcq}
  There is $\kappa_{\mathrm m}>0$ such that at every $x\in\Lset$ with $\theta(x)>0$
  there is a unit direction $d$ with
  $\langle\nabla c_i(x),d\rangle\le-\kappa_{\mathrm m}$ for all $i\in\A^+(x)$.
\end{assumption}

Assumption~\ref{ass:emfcq} is the natural extension of MFCQ
(Assumption~\ref{ass:mfcq}) from the active set at feasible limit points to the
positively-active set at infeasible points; at a feasible point ($\theta=0$,
$\A^+=\A$) it \emph{is} MFCQ. It is exactly the ``no infeasible stationary
points of $\theta$'' hypothesis that the feasible-start scoping
(Assumption~\ref{ass:feasstart}) presumes: it guarantees the funnel can always
make feasibility progress, so the method cannot stall at an infeasible point.

\begin{lemma}[MFCQ excludes infeasible $\theta$-stationarity]
  \label{lem:cq-nonstat}
  Under Assumption~\ref{ass:emfcq}, at every iteration the model feasibility
  criticality obeys
  \begin{equation}
    \label{eq:cq-nonstat}
    \chi_k^c \;=\; \|\nabla\theta_k^m(x_k)\|
             \;\ge\; \kappa_{\mathrm m}\sqrt{2\theta_k}\;-\;\kappa_{\mathrm{eg}}^\theta\Delta_k .
  \end{equation}
  In particular, for any $\varrho>0$, if $\theta_k\ge\varrho$ and
  $\Delta_k\le\bar\Delta(\varrho):=\kappa_{\mathrm m}\sqrt{2\varrho}/(2\kappa_{\mathrm{eg}}^\theta)$,
  then $\chi_k^c\ge\tfrac12\kappa_{\mathrm m}\sqrt{2\varrho}>0$.
\end{lemma}
\begin{proof}
  Because $\theta$ is smooth with $\nabla\theta(x_k)=J(x_k)^\top[c_k]_+=\sum_{i:c_i(x_k)>0}c_i(x_k)\nabla c_i(x_k)$,
  the extended-MFCQ direction $d$ (unit, $\langle\nabla c_i(x_k),d\rangle\le-\kappa_{\mathrm m}$
  for $i\in\A^+(x_k)\supseteq\{i:c_i(x_k)>0\}$) gives
  \[
    \|\nabla\theta(x_k)\| \;\ge\; -\langle\nabla\theta(x_k),d\rangle
      = -\sum_{i:c_i(x_k)>0}c_i(x_k)\langle\nabla c_i(x_k),d\rangle
      \;\ge\; \kappa_{\mathrm m}\!\!\sum_{i:c_i(x_k)>0}\!\!c_i(x_k)
      = \kappa_{\mathrm m}\|[c_k]_+\|_1 .
  \]
  Since $\|[c_k]_+\|_1\ge\|[c_k]_+\|_2=\sqrt{2\theta_k}$, we get
  $\|\nabla\theta(x_k)\|\ge\kappa_{\mathrm m}\sqrt{2\theta_k}$. The model gradient bound
  of L1, $\|\nabla\theta_k^m(x_k)-\nabla\theta(x_k)\|\le\kappa_{\mathrm{eg}}^\theta\Delta_k$,
  then yields \eqref{eq:cq-nonstat}. The ``in particular'' clause substitutes
  $\theta_k\ge\varrho$ and $\Delta_k\le\bar\Delta(\varrho)$.
\end{proof}

\noindent With the smooth measure this is a clean gradient bound --- no subgradient
or spurious-active-constraint bookkeeping. It is the exact model analog of the
Gould--Toint feasibility bound $\|\nabla\theta\|\ge\kappa_{\mathrm m}\sqrt{2\theta}$,
and it is the pivot that excludes infeasible $\theta$-stationary limit points in T1.

We give a \emph{self-contained} proof that $\theta_k\to0$, replacing the imported
dichotomy of the previous draft. We follow the logical order of
\citet{GouldToint2010} --- optimality progress is analyzed \emph{before}
feasibility, so no circular dependence on T2 arises --- adapt every step to the
derivative-free ($f,c$ known only through fully-linear models, via L1) and smooth
inequality ($\theta=\tfrac12\|[c]_+\|^2$) setting, and use only the reproduced
trust-region lemmas of Section~\ref{sec:tr-lemmas} (nothing imported). Write
$\kappa_{\mathrm{dec}}=\kappa_{\mathrm{tg}}\kappa_{\mathrm{tc}}$ for the
$f$-iteration decrease constant. On a successful $f$-iteration $j$, the ratio test
gives $f(x_j)-f(x_{j+1})\ge\eta_1\pred^f_j$; the $f$-domination test
\eqref{eq:f-domination} gives $\pred^f_j\ge\kappa_{\mathrm{tg}}\delta^{f,t}_j$ where
$\delta^{f,t}_j=m_f^j(x_j+n_j)-m_f^j(x_j+s_j)$ is the tangential objective decrease;
and L3 bounds $\delta^{f,t}_j$ below. Chaining,
\begin{equation}
  \label{eq:f-decrease}
  f(x_j)-f(x_{j+1}) \;\ge\; \eta_1\kappa_{\mathrm{tg}}\,\delta^{f,t}_j
    \;\ge\; \eta_1\kappa_{\mathrm{dec}}\,\pi_j^f\min\!\Bigl(\tfrac{\pi_j^f}{1+\kappa_H},\Delta_j\Bigr).
\end{equation}
(Comparing $\pred^f_j$ to $\delta^{f,t}_j$ --- not to $\pred^\theta_j$ --- is what
makes this bound in terms of $\pi_j^f$ valid; see \eqref{eq:f-domination}.)

\begin{lemma}[L6 --- path length between criticality crossings]
  \label{lem:pathlen}
  Let $k_1<k_2$ be successful iterations such that every successful iteration in
  $[k_1,k_2)$ is an $f$-iteration with $\pi_j^f\ge\varepsilon>0$, and
  $f(x_{k_1})-f(x_{k_2})\le\eta_1\kappa_{\mathrm{dec}}\varepsilon^2/\bigl(2(1+\kappa_H)\bigr)$.
  Then
  \begin{equation}
    \label{eq:pathlen}
    \|x_{k_1}-x_{k_2}\| \;\le\; \frac{1}{\eta_1\kappa_{\mathrm{dec}}\varepsilon}\bigl(f(x_{k_1})-f(x_{k_2})\bigr).
  \end{equation}
\end{lemma}
\begin{proof}
  For a successful $f$-iteration $j\in[k_1,k_2)$, \eqref{eq:f-decrease} and
  $\pi_j^f\ge\varepsilon$ give
  $f(x_j)-f(x_{j+1})\ge\eta_1\kappa_{\mathrm{dec}}\varepsilon\min(\varepsilon/(1+\kappa_H),\Delta_j)$.
  Were the minimum ever $\varepsilon/(1+\kappa_H)$, that single term would be
  $\ge\eta_1\kappa_{\mathrm{dec}}\varepsilon^2/(1+\kappa_H)=2\bigl[\eta_1\kappa_{\mathrm{dec}}\varepsilon^2/(2(1+\kappa_H))\bigr]\ge2\bigl(f(x_{k_1})-f(x_{k_2})\bigr)$,
  exceeding the total decrease $\sum_j(f(x_j)-f(x_{j+1}))=f(x_{k_1})-f(x_{k_2})$ (all
  terms nonnegative) --- impossible. So the minimum is always $\Delta_j$, whence
  $\|x_j-x_{j+1}\|=\|s_j\|\le\Delta_j\le(f(x_j)-f(x_{j+1}))/(\eta_1\kappa_{\mathrm{dec}}\varepsilon)$.
  On unsuccessful and $\mu$-iterations $x$ is unchanged. Summing over $j\in[k_1,k_2)$
  and using the triangle inequality gives \eqref{eq:pathlen}.
\end{proof}
This is the derivative-free, inequality analog of \citet[Lemma~3.20]{GouldToint2010}.

\begin{lemma}[L7 --- feasibility when finitely many successful $c$-iterations]
  \label{lem:feas-B}
  If $|\S_c|<\infty$ (Regime~B), then $\theta_k\to0$.
\end{lemma}
\begin{proof}
  Let $k_0$ be the last successful $c$-iteration. For $k>k_0$ every successful
  iteration is an $f$-iteration, so $\{f(x_k)\}_{k>k_0}$ is non-increasing (only
  $f$-iterations move $x$, and they decrease $f$) and bounded below
  (Proposition~\ref{prop:bounded}); it converges to some $f_*$. By (TR2) applied to
  these $f$-iterations (criticality $\pi_k^f$, merit $f$),
  $\liminf_{k}\pi_k^f=0$.

  Suppose $\theta_k\not\to0$: there are $\varepsilon_0>0$ and infinitely many
  successful $f$-iterations $k>k_0$ with $\theta_k\ge\varepsilon_0$. On an
  $f$-iteration the switching test \eqref{eq:switch} held, so
  $\pi_k^f>\omega_{\mathrm t}(\theta_k)\ge\omega_{\mathrm t}(\varepsilon_0)>0$
  ($\omega_{\mathrm t}$ non-decreasing, $\omega_{\mathrm t}(t)>0$ for $t>0$). Fix
  $\varepsilon$ with
  $0<\varepsilon<\tfrac12\omega_{\mathrm t}(\varepsilon_0)$ and small enough that
  $\omega_{\mathrm t}^{-1}(\varepsilon)<\varepsilon_0/2$ (possible as
  $\omega_{\mathrm t}$ is strictly increasing near $0$). Choose $k_1>k_0$ with
  $\theta_{k_1}\ge\varepsilon_0$ and
  \begin{equation}
    \label{eq:B-choice}
    f(x_{k_1})-f_* \;\le\; \eta_1\kappa_{\mathrm{dec}}\varepsilon\,
       \min\!\Bigl(\tfrac{\varepsilon}{2(1+\kappa_H)},\ \tfrac{\varepsilon_0}{2L_\theta}\Bigr),
  \end{equation}
  where $L_\theta=\sup_{\Lset}\|\nabla\theta\|<\infty$ is the Lipschitz constant of the
  \emph{smooth} $\theta$; possible since $f(x_k)\downarrow f_*$. Let $k_2>k_1$ be the
  first successful iteration with $\pi_{k_2}^f<\varepsilon$ (it exists since
  $\liminf\pi_k^f=0$ and $\pi_{k_1}^f\ge\omega_{\mathrm t}(\varepsilon_0)>\varepsilon$);
  every successful iteration in $[k_1,k_2)$ is an $f$-iteration with $\pi_j^f\ge\varepsilon$.
  The first argument of \eqref{eq:B-choice} meets L6's hypothesis, so L6 gives
  $P:=\sum_{j\in[k_1,k_2)\cap\S}\|x_j-x_{j+1}\|\le(f(x_{k_1})-f_*)/(\eta_1\kappa_{\mathrm{dec}}\varepsilon)\le\varepsilon_0/(2L_\theta)$.
  Each accepted step segment $[x_j,x_{j+1}]$ lies in $B(x_j,\Delta_j)\subseteq\Lset$,
  where $\|\nabla\theta\|\le L_\theta$; the mean-value theorem on each segment and
  summation give
  \[
    \bigl|\theta_{k_1}-\theta_{k_2}\bigr|
      \le\sum_j|\theta(x_j)-\theta(x_{j+1})|\le L_\theta\sum_j\|x_j-x_{j+1}\|=L_\theta P\le\varepsilon_0/2,
  \]
  so $\theta_{k_2}\ge\theta_{k_1}-\varepsilon_0/2\ge\varepsilon_0/2$ (the polygonal path
  avoids assuming $\Lset$ convex).
  But $k_2$ is a successful $f$-iteration, so
  $\omega_{\mathrm t}(\theta_{k_2})<\pi_{k_2}^f<\varepsilon$, giving
  $\theta_{k_2}<\omega_{\mathrm t}^{-1}(\varepsilon)<\varepsilon_0/2$ --- a
  contradiction. Hence $\theta_k\to0$.
\end{proof}
This is the derivative-free, inequality analog of \citet[Lemma~3.21]{GouldToint2010};
the inequality enters only through the $1$-Lipschitz bound on $[\cdot]_+$.

\begin{lemma}[L8 --- feasibility when infinitely many successful $c$-iterations]
  \label{lem:feas-A}
  If $|\S_c|=\infty$ (Regime~A), then $\theta_k\to0$.
\end{lemma}
\begin{proof}
  By L4 the funnel converges, $\theta_k^{\max}\downarrow\theta_*^{\max}\ge0$. Suppose
  $\theta_*^{\max}=\tau>0$. On a successful $c$-iteration, L4-gap \eqref{eq:L4-gap}
  and $\theta_k\le\theta_k^{\max}$ telescope:
  \[
    \sum_{k\in\S_c}\ared^\theta_k
      \le\sum_{k\in\S_c}\bigl(\theta_k^{\max}-\theta_{k+1}\bigr)
      \le\frac{1}{1-\kappa_{\mathrm{cb}}}\bigl(\theta_0^{\max}-\tau\bigr)<\infty,
  \]
  so $\ared^\theta_k\to0$, hence $\pred^\theta_k\le\ared^\theta_k/\eta_1\to0$ on
  $\S_c$. Moreover, because $\kappa_{\mathrm{ca}}\theta_k^{\max}\to\kappa_{\mathrm{ca}}\tau<\tau$
  while $\theta_{k+1}^{\max}\ge\tau$, for large $k\in\S_c$ the maximum in
  \eqref{eq:funnel-update} is attained by its second argument, giving
  $\theta_{k+1}\ge(\tau-\kappa_{\mathrm{cb}}\theta_k^{\max})/(1-\kappa_{\mathrm{cb}})\to\tau$;
  thus $\theta_k\ge\theta_{k+1}\ge\tau/2$ for all large $k\in\S_c$.

  For all large $k\in\S_c$ we thus have $\theta_k\ge\tau/2>0$. A $c$-iteration
  computes a normal step, so the feasibility radius cap \eqref{eq:feas-cap} is in
  force, and (Section~\ref{sec:algo-crit}) it makes cq-nonstat spurious-free:
  \begin{equation}
    \label{eq:chi-tau}
    \chi_k^c\;\ge\;\frac{\kappa_{\mathrm m}\sqrt{2\theta_k}}{1+\mu\kappa_{\mathrm{eg}}^\theta}
      \;\ge\;\varsigma_\tau:=\frac{\kappa_{\mathrm m}\sqrt{\tau}}{1+\mu\kappa_{\mathrm{eg}}^\theta}>0
    \qquad(k\in\S_c\text{ large}),
  \end{equation}
  \emph{with no large-radius exception} --- the cap $\Delta_k\le\mu\chi_k^c$ is exactly
  what prevents an over-large radius from masking the true criticality (this is where
  the smooth measure and the feasibility criticality step earn their keep). Every
  $c$-iteration carries a genuine normal step (Section~\ref{sec:algo-accept}), so the
  normal Cauchy bound \eqref{eq:normal-cauchy} and feasibility-retention
  \eqref{eq:retention} ($\pred^\theta_k\ge(1-\kappa_{\mathrm{nt}})\times$\eqref{eq:normal-cauchy}) give
  \begin{equation}
    \label{eq:A-predtheta}
    \pred^\theta_k\;\ge\;\tfrac{1-\kappa_{\mathrm{nt}}}{2}\,\chi_k^c\min\!\Bigl(\tfrac{\chi_k^c}{1+\beta_\theta},\kappa_{\mathrm n}\Delta_k\Bigr)
    \qquad\text{with } \chi_k^c\ge\varsigma_\tau .
  \end{equation}

  \emph{Step 1 (summable radii on $\S_c$).} Since $\chi_k^c\ge\varsigma_\tau$ is bounded
  below on large $\S_c$, \eqref{eq:A-predtheta} gives
  $\pred^\theta_k\ge\tfrac{1-\kappa_{\mathrm{nt}}}{2}\varsigma_\tau\min(\varsigma_\tau/(1+\beta_\theta),\kappa_{\mathrm n}\Delta_k)$.
  At most finitely many $k\in\S_c$ have $\Delta_k\ge\varsigma_\tau/(\kappa_{\mathrm n}(1+\beta_\theta))$
  (each such has $\pred^\theta_k$ bounded below, but $\pred^\theta_k\to0$); for the rest
  the $\min$ is $\kappa_{\mathrm n}\Delta_k$, so $\pred^\theta_k\ge c_2\Delta_k$ with
  $c_2:=\tfrac{(1-\kappa_{\mathrm{nt}})\kappa_{\mathrm n}\varsigma_\tau}{2}$. With the
  telescoping bound,
  $\sum_{k\in\S_c}\Delta_k\le c_2^{-1}\sum_{k\in\S_c}\pred^\theta_k+O(1)<\infty$.
  \emph{(No trust-region radius-bookkeeping is imported here: the cap
  \eqref{eq:feas-cap} makes $\chi_k^c$ genuinely bounded below, so the earlier
  ``large-radius spurious criticality'' case cannot arise.)}\hfill$(\ast)$

  \emph{Step 2 (finitely many band up-crossings).} Let $L_f=\sup_{\Lset}\|\nabla f\|$
  and $L_\theta=\sup_{\Lset}\|\nabla\theta\|$ (finite, Assumptions~\ref{ass:smooth},
  \ref{ass:levelset}), $\pi_B:=\omega_{\mathrm t}(\tau/4)>0$, and consider the band
  $\{x:\tau/4\le\theta(x)\le\tau/2\}$. An \emph{up-crossing} is an ascent from
  $\theta\le\tau/4$ to $\theta\ge\tau/2$. Only $f$-iterations raise $\theta$ (a
  $c$-iteration reduces it; $\mu$/unsuccessful leave $x$ fixed), and a \emph{rising}
  $f$-iteration in the band has, by switching \eqref{eq:switch},
  $\pi_j^f>\omega_{\mathrm t}(\theta_j)\ge\pi_B$; hence by \eqref{eq:f-decrease}, when
  $\Delta_j\le\pi_B/(1+\kappa_H)$,
  \[
    f(x_j)-f(x_{j+1})\ge\eta_1\kappa_{\mathrm{dec}}\pi_B\Delta_j
      \;\ge\;\tfrac{\eta_1\kappa_{\mathrm{dec}}\pi_B}{L_\theta}\,\bigl(\theta_{j+1}-\theta_j\bigr)
  \]
  (using $\theta_{j+1}-\theta_j\le L_\theta\|s_j\|\le L_\theta\Delta_j$), and
  $f(x_j)-f(x_{j+1})\ge\eta_1\kappa_{\mathrm{dec}}\pi_B^2/(1+\kappa_H)$ otherwise.
  Summing the rising $f$-iterations of one up-crossing (whose $\theta$-increases total
  $\ge\tau/4$) yields an objective decrease $\ge
  \Phi:=\min\!\bigl(\tfrac{\eta_1\kappa_{\mathrm{dec}}\pi_B}{L_\theta}\tfrac{\tau}{4},\ \tfrac{\eta_1\kappa_{\mathrm{dec}}\pi_B^2}{1+\kappa_H}\bigr)>0$
  per up-crossing. Meanwhile $f$ \emph{increases} only on successful $c$-iterations,
  each by $\le L_f\Delta_k$, so the total increase is $\le L_f\sum_{k\in\S_c}\Delta_k<\infty$
  by $(\ast)$. Since $f\ge f_{\mathrm{low}}$ (Proposition~\ref{prop:bounded}),
  \[
    (\#\text{up-crossings})\cdot\Phi \le f(x_0)-f_{\mathrm{low}}+L_f\textstyle\sum_{k\in\S_c}\Delta_k <\infty,
  \]
  so there are finitely many up-crossings.

  \emph{Step 3 (contradiction).} After the last up-crossing, either
  \emph{(I)} $\theta_k\ge\tau/4$ for all large $k$, or \emph{(II)} $\theta_k<\tau/2$
  for all large $k$.

  In case~(II) only finitely many $k$ have $\theta_k\ge\tau/2$; but every large
  $k\in\S_c$ has $\theta_k\ge\tau/2$, so $|\S_c|<\infty$ --- contradicting Regime~A.

  In case~(I) all large iterations lie in the \emph{pure} infeasible region
  $\{\theta\ge\tau/4\}$ (no near-feasible excursions interleave). There
  cq-nonstat \eqref{eq:cq-nonstat} gives, for $\Delta_k\le\varrho_0:=\kappa_{\mathrm m}\sqrt{\tau/2}/(2\kappa_{\mathrm{eg}}^\theta)$,
  $\chi_k^c\ge\tfrac12\kappa_{\mathrm m}\sqrt{\tau/2}=:\varsigma_0>0$; so no $\mu$-iteration
  occurs (Lemma~\ref{lem:mu-critical} needs $\chi_k^c=0$), and $\Delta$ can shrink only
  on \emph{unsuccessful} iterations. Apply the reproduced accuracy lemma
  \ref{lem:tr-acc}: a $c$-iteration (merit $\theta$, $\xi=\chi_k^c\ge\varsigma_0$,
  decrease \eqref{eq:A-predtheta}) and an $f$-iteration (merit $f$, $\xi=\pi_k^f>\pi_B$
  by switching, decrease \eqref{eq:f-decrease}) are both \emph{successful} once
  $\Delta_k\le\Delta_3:=\min(\varrho_0,\bar\Delta(\varsigma_0),\bar\Delta(\pi_B))$
  (the accuracy thresholds of Lemma~\ref{lem:tr-acc} for the two merits). Hence no
  shrinkage occurs at $\Delta_k\le\Delta_3$, so by Lemma~\ref{lem:tr1}
  $\Delta_k\ge\gamma\Delta_3=:\Delta_{\mathrm{lb}}>0$ for all large $k$. But this
  contradicts $(\ast)$ ($\sum_{k\in\S_c}\Delta_k<\infty$, with $|\S_c|=\infty$, forces
  $\Delta_k\to0$ along $\S_c$).

  Both cases are impossible, so $\tau=0$: $\theta_k^{\max}\to0$ and
  $\theta_k\le\theta_k^{\max}\to0$.
\end{proof}

\begin{theorem}[T1 --- feasibility]
  \label{thm:T1}
  Under Assumptions~\ref{ass:relaxable}--\ref{ass:mfcq} and \ref{ass:emfcq},
  $\theta_k\to0$; hence $\theta_*^{\max}=0$ and every limit point $x_*$ of
  $\{x_k\}$ is feasible, $c(x_*)\le0$.
\end{theorem}
\begin{proof}
  Either $|\S_c|<\infty$ (L7) or $|\S_c|=\infty$ (L8); both give $\theta_k\to0$. By
  continuity of $c$, $c(x_*)=\lim c(x_k)\le0$ at every limit point.
\end{proof}

\begin{remark}[The feasibility theorem is self-contained]
  Theorem~\ref{thm:T1} is proved in full. Every trust-region ingredient it uses ---
  the accuracy-implies-success lemma \ref{lem:tr-acc}, the radius lower bound
  \ref{lem:tr1}, and the liminf-criticality \ref{lem:tr2} (invoked in Regime~B) --- is
  \emph{reproduced} in Section~\ref{sec:tr-lemmas}; nothing is imported. The smooth
  measure $\theta=\tfrac12\|[c]_+\|^2$ is what makes this clean: it turns feasibility
  into smooth trust-region minimization with the gradient criticality
  $\chi^c=\|\nabla\theta^m\|$, and the feasibility radius cap \eqref{eq:feas-cap}
  removes the large-radius ``spurious criticality'' obstruction outright (the former
  residual TD1$'$ is gone --- see Step~1 of L8). The genuinely new content is the
  extended-MFCQ pivot (Lemma~\ref{lem:cq-nonstat}, now a one-line gradient bound
  $\chi^c\ge\kappa_{\mathrm m}\sqrt{2\theta}-\kappa_{\mathrm{eg}}^\theta\Delta$), the
  funnel--inequality interaction, and the band-crossing argument that tames the
  interleaving of infeasible and near-feasible iterations.

  The pivot throughout is Lemma~\ref{lem:cq-nonstat}: extended MFCQ forces
  $\chi_k^c>0$ at infeasible points, ruling out the ``infeasible stall'' a pure
  funnel-relative trigger would permit. This is where the feasible-reachable-start
  scoping (Assumption~\ref{ass:feasstart}) pays off: it is consistent with the absence
  of infeasible $\theta$-stationary traps that Assumption~\ref{ass:emfcq} formalizes.
\end{remark}

\subsection{Optimality: a KKT limit point (Theorem T2)}
\label{sec:T2}

With feasibility secured, the optimality argument applies the reproduced
trust-region lemmas of Section~\ref{sec:tr-lemmas} to the $f$-iterations, which
dominate once $\theta_k\to0$.

\begin{lemma}[Eventually $f$-iterations near a non-KKT accumulation]
  \label{lem:eventually-f}
  If $\pi_k^f\ge\pi_*>0$ along a subsequence, then along that subsequence the
  switching condition \eqref{eq:switch} holds for all large $k$ and each such
  iteration is an $f$-iteration.
\end{lemma}
\begin{proof}
  By T1, $\theta_k\to0$. Since $\pi_k^f\ge\pi_*$ and $\omega_{\mathrm t}(\theta_k)\to
  \omega_{\mathrm t}(0)=0$, the switching test \eqref{eq:switch}
  $\pi_k^f>\omega_{\mathrm t}(\theta_k)$ holds for all large $k$ in the subsequence, so
  a tangential step is taken and, by L3,
  \[
    \delta^{f,t}_k \;\ge\; \kappa_{\mathrm{tc}}\,\pi_*\min\!\bigl(\pi_*/(1+\kappa_H),\Delta_k\bigr)\;>\;0.
  \]
  We must show the $f$-domination test \eqref{eq:f-domination},
  $\pred^f_k\ge\kappa_{\mathrm{tg}}\delta^{f,t}_k$, so the iteration is classified $f$.
  Writing $\pred^f_k=\delta^{f,t}_k+[m_f^k(x_k)-m_f^k(x_k+n_k)]$, the normal-step
  objective change is bounded using $\|n_k\|\le\kappa_{\mathrm{nn}}\sqrt{2\theta_k}$
  \eqref{eq:normal-size} and $\|g_k\|,\|H_k\|$ bounded (Proposition~\ref{prop:bounded}):
  \[
    \bigl|m_f^k(x_k)-m_f^k(x_k+n_k)\bigr|
      \le \|g_k\|\,\|n_k\| + \tfrac12\|H_k\|\,\|n_k\|^2
      \le C\,\sqrt{\theta_k}
  \]
  for a constant $C$ (as $\theta_k\to0$, $\sqrt{\theta_k}$ dominates $\theta_k$). Near a
  non-critical accumulation the criticality step (L5) keeps
  $\Delta_k\ge\kappa_\Delta\sigma_k\ge\kappa_\Delta\pi_*$ bounded below, so
  $\delta^{f,t}_k\ge\kappa_{\mathrm{tc}}\pi_*\min(\pi_*/(1+\kappa_H),\kappa_\Delta\pi_*)=:\delta_*>0$
  is bounded below, while $C\sqrt{\theta_k}\to0$ (by T1). Hence for all large $k$,
  $C\sqrt{\theta_k}\le(1-\kappa_{\mathrm{tg}})\delta_*\le(1-\kappa_{\mathrm{tg}})\delta^{f,t}_k$,
  giving $\pred^f_k\ge\delta^{f,t}_k-C\sqrt{\theta_k}\ge\kappa_{\mathrm{tg}}\delta^{f,t}_k$.
  Thus \eqref{eq:f-domination} holds and the iteration is an $f$-iteration.
\end{proof}

\begin{theorem}[T2 --- first-order optimality / KKT]
  \label{thm:T2}
  Under Assumptions~\ref{ass:relaxable}--\ref{ass:mfcq} and \ref{ass:emfcq},
  \[
    \liminf_{k\to\infty}\ \pi_k^f \;=\; 0 .
  \]
  Consequently there is a limit point $x_*$ of $\{x_k\}$ that is feasible
  ($c(x_*)\le0$, by T1) and satisfies the KKT conditions of \eqref{eq:problem}: there
  exists $\lambda\ge0$ with $\nabla f(x_*)+\sum_i\lambda_i\nabla c_i(x_*)=0$ and
  $\lambda_i c_i(x_*)=0$. Under the criticality step (Lemma~\ref{lem:L5}) the stronger
  $\lim_{k\to\infty}\pi_k^f=0$ holds, so \emph{every} feasible limit point is KKT.
\end{theorem}
\begin{proof}
  \emph{$\liminf\pi_k^f=0$.} Suppose not: $\pi_k^f\ge\pi_*>0$ for all large $k$.
  By Lemma~\ref{lem:eventually-f} all large iterations are $f$-iterations, so $f$ is
  non-increasing over the tail (only $f$-iterations move $x$, decreasing $f$; the
  others fix it) and bounded below (Proposition~\ref{prop:bounded}); and L3 supplies
  the fraction-of-Cauchy decrease \eqref{eq:fcd} for the merit $f$ with criticality
  $\pi^f$. The reproduced Lemma~\ref{lem:tr2} (TR2), applied with $\varphi=f$,
  $\xi=\pi^f$, then gives $\liminf_k\pi_k^f=0$ --- contradicting $\pi_k^f\ge\pi_*$.

  \emph{KKT limit point.} Take a subsequence with $\pi_k^f\to0$ and (by boundedness
  of $\Lset$) $x_k\to x_*$. By T1, $\theta(x_*)=0$, so $x_*$ is feasible. On this
  subsequence the criticality step drives $\Delta_k\to0$ (L5, since
  $\sigma_k\to0$), so L1 gives $g_k\to\nabla f(x_*)$ and $J_k\to\nabla c(x_*)$.

  \emph{Active-set identification (shrinking $\varepsilon_k$).} The model
  active set uses the \emph{shrinking} tolerance
  $\varepsilon_k=\kappa_{\eps}\sigma_k\to0$ (Section~\ref{sec:algo-data}), so
  $\A_k(\varepsilon_k)=\{i:c_i(x_k)\ge-\varepsilon_k\}$. For $i\notin\A(x_*)$ we have
  $c_i(x_*)<0$; by continuity $c_i(x_k)\to c_i(x_*)<0$, so for large $k$,
  $c_i(x_k)<-\varepsilon_k$ and $i\notin\A_k(\varepsilon_k)$. Conversely each
  $i\in\A(x_*)$ ($c_i(x_*)=0$) satisfies $c_i(x_k)\ge-\varepsilon_k$ for large $k$.
  Hence $\A_k(\varepsilon_k)=\A(x_*)$ for all large $k$ in the subsequence --- exact
  identification, which a \emph{fixed} $\varepsilon$ would not deliver.

  \emph{Passing to the limit.} The model KKT residual $\pi_k^f\to0$ furnishes
  multipliers $\lambda^k\ge0$ supported on $\A_k(\varepsilon_k)=\A(x_*)$ with
  $g_k+(J_k^{\A})^\top\lambda^k\to0$. Since $\A_k(\varepsilon_k)=\A(x_*)$ and
  $J_k^{\A}\to\nabla c_{\A}(x_*)$, MFCQ (Assumption~\ref{ass:mfcq}) at $x_*$ ---
  linear independence of a strictly-feasible direction, equivalently boundedness of
  the multiplier set --- bounds $\{\lambda^k\}$, so it has a limit point
  $\lambda\ge0$ supported on $\A(x_*)$. Passing to the limit in
  $g_k+(J_k^{\A})^\top\lambda^k\to0$ gives $\nabla f(x_*)+\nabla c(x_*)^\top\lambda=0$;
  complementarity $\lambda_ic_i(x_*)=0$ holds because $\lambda_i=0$ for
  $i\notin\A(x_*)$ (multipliers supported on $\A(x_*)$) and $c_i(x_*)=0$ for
  $i\in\A(x_*)$. Thus $x_*$ is a KKT point.

  \emph{Strengthening to $\lim\pi_k^f=0$.} Once $\theta_k\to0$ the iterations are
  eventually $f$-iterations wherever $\pi_k^f$ is bounded below
  (Lemma~\ref{lem:eventually-f}), so the reproduced Lemma~\ref{lem:tr-lim} (TR-lim)
  applies with $\varphi=f$, $\xi=\pi^f$, giving $\lim_k\pi_k^f=0$; hence \emph{every}
  limit point is a KKT point.
\end{proof}

\subsection{Limiting-case validations}
\label{sec:limits}

We check that the analysis reduces correctly in the four regimes named in the
plan; these are the sanity constraints on the proof.

\paragraph{(a) $p=0$ (unconstrained) $\Rightarrow$ CSV Ch.~10/11.}
With no constraints, $\theta\equiv0$, $\chi_k^c\equiv0$, $n_k\equiv0$,
$\theta_k^{\max}$ is irrelevant, and every iteration is an $f$-iteration with
$\pi_k^f=\|g_k\|$ and $\A_k=\varnothing$. Algorithm~1 becomes
\path{Algorithm_11_1.m}, L3 becomes the unconstrained Cauchy decrease, L5 becomes
the CSV criticality step, T1 is vacuous, and T2 reduces to
$\liminf\|g_k\|=0$ (resp.\ $\lim\|g_k\|=0$), i.e.\ CSV Ch.~10/11.

\paragraph{(b) Equalities $\Rightarrow$ Sampaio--Toint (2015).}
An equality $c^{\mathrm{eq}}(x)=0$ is $\{c^{\mathrm{eq}}\le0,\,-c^{\mathrm{eq}}\le0\}$;
then $[c]_+$ recovers $|c^{\mathrm{eq}}|$ componentwise, so
$\theta=\tfrac12\|[c]_+\|^2=\tfrac12\|c^{\mathrm{eq}}\|^2$ with
$\nabla\theta=(J^{\mathrm{eq}})^\top c^{\mathrm{eq}}$ and feasibility criticality
$\chi^c=\|(J^{\mathrm{eq}})^\top c^{\mathrm{eq}}\|$. This is \emph{exactly} the squared
infeasibility measure $\tfrac12\|c\|^2$ and criticality $\|J^\top c\|$ of
\citet{SampaioToint2015}: having adopted the smooth measure, the correspondence is now
a genuine specialization (not merely structural), with L2/L4/T1 specializing to their
feasibility lemmas and T2 to their first-order result. The only differences are the
inequality one-sidedness ($[\cdot]_+$) and the derivative-free model layer (L1).

\paragraph{(c) MFCQ failure $\Rightarrow$ identifies the breaking lemma.}
If Assumption~\ref{ass:emfcq} fails at some infeasible $\bar x\in\Lset$
(no direction strictly reduces all positively-active constraints), then
Lemma~\ref{lem:cq-nonstat} fails: $\chi_k^c$ may vanish with $\theta_k>0$, so L2's
right-hand side is $0$ and T1's feasibility drive stalls --- the algorithm may
converge to an \emph{infeasible} stationary point of $\theta$. This is precisely
the dropped dichotomy branch; it localizes the indispensability of MFCQ to
Lemma~\ref{lem:cq-nonstat} $\to$ T1. If instead MFCQ (Assumption~\ref{ass:mfcq})
fails only at a \emph{feasible} limit point, T1 still holds but the multiplier set
in T2 may be unbounded, breaking the KKT conclusion (only Fritz--John stationarity
survives).

\paragraph{(d) Feasible-iterate regime $\Rightarrow$ NOWPAC relation.}
If the acceptance rule is tightened to admit only $\theta(x_k+s_k)=0$ (equivalently
$\theta_k^{\max}\equiv0$), then no relaxable excursion is ever taken, the normal
step keeps iterates feasible, and Algorithm~1 runs entirely inside $\feas$. This is
the feasible-iterate regime of NOWPAC \citep{AugustinMarzouk2017,Menhorn2022};
our analysis then does not \emph{use} relaxability, and T1 is automatic
($\theta_k\equiv0$). The contrast is that our default ($\theta_k^{\max}>0$)
\emph{exploits} relaxability to sample and step through infeasible points, which is
unavailable to a feasible-iterate method.

\subsection{Summary of the argument}
\label{sec:summary}

The charter question --- ``do fully-linear models of $f$ and $c$ on every iteration
suffice for convergence to a KKT point?'' --- is answered: \emph{necessary, but
only in concert with two further mechanisms}. L1 shows fully-linearity of $f$ and
$c$ (and the inherited accuracy of $\theta$) controls all four ratio tests; L5
(the criticality step) ties $\Delta_k$ to criticality so model stationarity
transfers to true stationarity; and the feasibility theorem~\ref{thm:T1} ---
proved self-contained via the path-length lemma~\ref{lem:pathlen}, the two
regime lemmas~\ref{lem:feas-B}--\ref{lem:feas-A}, and the extended-MFCQ
Lemma~\ref{lem:cq-nonstat} --- drives $\theta_k\to0$. Remove any one and a
limiting-case check in Section~\ref{sec:limits} breaks. The two theorems T1
(feasibility) and T2 (KKT under MFCQ) are the deliverable D3. \emph{The analysis is
now self-contained}: the single-region trust-region facts it needs
(Lemmas~\ref{lem:tr-acc}--\ref{lem:tr-lim}) are reproduced in
Section~\ref{sec:tr-lemmas} from first principles, and the smooth measure
$\theta=\tfrac12\|[c]_+\|^2$ with the feasibility radius cap \eqref{eq:feas-cap}
removes the large-radius obstruction that the earlier unsquared draft could only cite.


\section*{Status}
\addcontentsline{toc}{section}{Status}
Milestones M1--M4 are complete (problem/notation, algorithm, convergence analysis,
and an adversarial self-review with the D5 code-bridge appendix), and both prior
technical-debt items are now closed:

\textbf{TD1 (self-contained feasibility) --- resolved.} The feasibility
theorem~\ref{thm:T1} ($\theta_k\to0$) is proved directly in both regimes
(Section~\ref{sec:T1}): Lemma~\ref{lem:pathlen} (path length), \ref{lem:feas-B}
(finitely many successful $c$-iterations), \ref{lem:feas-A} (infinitely many), in the
Gould--Toint logical order (optimality before feasibility). The \emph{interleaving}
of infeasible and near-feasible iterations is handled self-contained by a
summable-radii bound and a \emph{finitely-many-band-crossings} argument; the imported
trust-funnel dichotomy is gone.

\textbf{TD2 (single-region trust-region machinery) --- resolved.} Two design changes
requested and made: (i) the \emph{smooth} infeasibility measure
$\theta=\tfrac12\|[c]_+\|^2$ (Section~\ref{sec:infeas}), which is $C^1$ with
$\nabla\theta=J^\top[c]_+$, turns feasibility into smooth trust-region minimization
with the gradient criticality $\chi^c=\|\nabla\theta^m\|$ and makes the MFCQ pivot
(Lemma~\ref{lem:cq-nonstat}) a one-line bound
$\chi^c\ge\kappa_{\mathrm m}\sqrt{2\theta}-\kappa_{\mathrm{eg}}^\theta\Delta$; and (ii)
the standard single-region trust-region results are now \emph{reproduced from first
principles} in Section~\ref{sec:tr-lemmas} (Lemmas~\ref{lem:tr-acc}
accuracy$\Rightarrow$success, \ref{lem:tr1} radius lower bound, \ref{lem:tr2} liminf
criticality, \ref{lem:tr-lim} full criticality limit), applied to both merits $f$ and
$\theta$. The feasibility radius cap \eqref{eq:feas-cap} removes the former
large-radius ``spurious criticality'' obstruction (the earlier TD1$'$) outright.

\noindent\textbf{The convergence analysis is now self-contained}, resting only on the
stated assumptions and results proved within: T1 (feasibility, $\theta_k\to0$) and T2
($\lim_k\pi_k^f=0$, every feasible limit point KKT under MFCQ). The
constrained/relaxable novelty --- extended MFCQ excluding infeasible
$\theta$-stationarity, the funnel--inequality interaction, and the band-crossing
feasibility argument --- is proved in full; the reproduced single-region lemmas cite
\citet[Ch.~6.4]{ConnGouldToint2000} only as the source whose arguments they restate.

% ---- D5 code-bridge appendix (M4) ----
% !TEX root = main.tex
%======================================================================
%  D5 --- Code-bridge appendix
%  Maps each Algorithm 1 component to existing IBCDFO building blocks.
%  THEORY PHASE ONLY: no solver code is written here; this is the
%  implementation map for the later coding phase.
%======================================================================
\appendix
\section{Code-bridge appendix (D5)}
\label{sec:codebridge}

This appendix maps each component of Algorithm~1 (Section~\ref{sec:algorithm}) to
an existing routine in this repository, so the later implementation phase is a
translation rather than a fresh build. \textbf{No code is written in this
theory phase.} Paths are relative to the repository root; all listed files were
confirmed present at the time of writing. Each row also flags whether the block is
\emph{reusable as-is}, needs \emph{adaptation}, or is a genuine \emph{extension
point} (new code the later phase must add).

\begin{center}\small
\renewcommand{\arraystretch}{1.25}
\begin{tabular}{@{}p{0.30\textwidth} p{0.42\textwidth} p{0.19\textwidth}@{}}
\toprule
\textbf{Algorithm~1 component} & \textbf{Existing building block} & \textbf{Status}\\
\midrule
Step~1 fully-linear models of $F=(f,c_1,\dots,c_p)$; $g_k,J_k$, model Hessians; \verb|valid| flag (Assumption~\ref{ass:fl})
  & \path{pounders/py/formquad.py} --- vector-valued minimum-Frobenius-norm models, returns \verb|G|, \verb|H|, and \verb|valid|; monomial basis in \path{phi2eval.py}
  & reuse as-is \\
$\Lambda$-poised geometry management (uniform constants of L1/L5)
  & \path{DFO_quad_TR/}: \path{Algorithm_6_4.m}, \path{Model_Improvement_Step.m}, \path{Maximize_Lagrange_polynomial.m}, \path{Test_fully_quadratic.m}; geometry rays via \path{pounders/py/bmpts.py} + \path{boxline.py}
  & reuse as-is \\
Infeasibility measure $\theta=\|[c]_+\|$ and its model $\theta_k^m$ \eqref{eq:theta}--\eqref{eq:theta-model}
  & compose from the vector model $m_c^k$ (from \path{formquad}); the composition pattern mirrors \path{pounders/py/general_h_funs.py} \verb|combine_*| (e.g.\ \verb|combine_leastsquares|)
  & adapt \\
Two criticality measures $\chi_k^c$ \eqref{eq:chi-c}, $\pi_k^f$ \eqref{eq:pi-f}
  & new: convex-composite steepest-descent slope for $\chi_k^c$; projected-gradient/KKT residual for $\pi_k^f$ over $J_k^{\A}$
  & extension point \\
Step~1 criticality loop (L5), $\Delta$ tie-in \eqref{eq:crit-radius}
  & \path{DFO_quad_TR/Algorithm_10_2.m} --- structure and $\Delta=\min(\max(\Delta,\beta\sigma),\Delta_{\text{in}})$ carried over verbatim; only $\sigma_k=\max(\chi_k^c,\pi_k^f)$ replaces the scalar criticality
  & adapt \\
Step~2a normal step $n_k$: $\min\|[c_k+J_k n]_+\|$ s.t.\ $\|n\|\le\kappa_{\mathrm n}\Delta$ \eqref{eq:normal-sub}
  & convex QP/LP in $(c_k,J_k)$; box-constrained core solvers \path{pounders/py/bqmin.py} (\verb|spsolver=1|) or external \path{minq} submodule \verb|minqsw| (\verb|spsolver=2|, Neumaier minq5); needs an inequality/violation layer on top
  & extension point \\
Step~2b tangential step $t_k$: reduce $m_f^k$ on $\{J_k^{\A}t\le0\}\cap B$ \eqref{eq:tang-sub}
  & quadratic-model minimization over box: \path{bqmin.py}/\verb|minqsw|; the null-space/active-cone restriction $J_k^{\A}t\le0$ is added on top
  & adapt \\
Step~3 ratios $\rho^f_k,\rho^c_k$; $f$-/$c$-iteration classification \eqref{eq:rhof}--\eqref{eq:rhoc}
  & pattern of \path{DFO_quad_TR/Algorithm_11_1.m} Step~3 (\verb|rho| test); add the $\theta$-ratio $\rho^c_k$ and the switching-based classification
  & adapt \\
Step~4b funnel update $\theta_{k+1}^{\max}$ \eqref{eq:funnel-update}
  & new scalar bookkeeping (Gould--Toint rule)
  & extension point \\
Step~5 radius update \eqref{eq:radius-update}
  & \path{DFO_quad_TR/Algorithm_11_1.m} Step~5 --- carried over verbatim
  & reuse as-is \\
Blackbox eval of $F$ (one call returns $f$ and all $c_i$)
  & driver pattern in \path{Algorithm_11_1.m}; input checks \path{pounders/py/checkinputss.py}
  & reuse as-is \\
\bottomrule
\end{tabular}
\end{center}

\paragraph{Deliberate non-reuse (penalty).}
The repository contains a penalty-style composition,
\verb|h_max_plus_quadratic_violation_penalty| in
\path{manifold_sampling/py/general_nonsmooth_h_funs.py}. Consistent with the
design decision recorded in Section~\ref{sec:positioning} (trust-funnel,
\emph{not} penalty), Algorithm~1 does \textbf{not} use it: infeasibility is
controlled by the funnel bound $\theta_k^{\max}$, never folded into the objective
through a penalty parameter. It is listed here only to mark it as explicitly out
of scope.

\paragraph{Manifold/active-set machinery for the nonsmooth $\theta$.}
The nonsmoothness of $\theta=\|[c]_+\|$ at active/violated boundaries
(Section~\ref{sec:infeas}) is handled in the analysis by the convex-composite
argument. For a robust implementation of $\chi_k^c$ and the model active set
$\A^+$, the manifold-sampling machinery
(\path{manifold_sampling/py/build_p_models.py}, \verb|h_pw_maximum| and relatives
in \path{general_nonsmooth_h_funs.py}) is the natural tool for identifying the
active manifold of $[\cdot]_+$ --- reused for its \emph{geometry/active-set}
identification, not for any penalty semantics.

\paragraph{Summary of extension points (the new code the next phase must write).}
\begin{enumerate}\itemsep1pt
  \item the two criticality measures $\chi_k^c,\pi_k^f$;
  \item the linearized-violation normal-step subproblem \eqref{eq:normal-sub} on
        top of the box-QP solvers;
  \item the active-cone restriction for the tangential subproblem \eqref{eq:tang-sub};
  \item the funnel bookkeeping (bound $\theta_k^{\max}$, update, and the
        $f$-/$c$-iteration switch).
\end{enumerate}
Everything else --- model building, poisedness/geometry, the criticality-loop
skeleton, ratio tests, and the radius update --- is already present and reused.


\bibliographystyle{plainnat}
\bibliography{refs}

\end{document}
